\documentclass[11pt]{article}
\usepackage[a4paper,margin=24mm]{geometry}
\usepackage[T1]{fontenc}
\usepackage{lmodern,amsmath,amssymb,amsthm,mathtools}
\usepackage{booktabs,longtable,array,enumitem,microtype,xcolor}
\usepackage{xurl,hyperref,listings,fancyhdr}
\hypersetup{colorlinks=true,linkcolor=blue!45!black,citecolor=blue!45!black,urlcolor=blue!45!black,pdftitle={Centered residue certificates and reusable Pell denominators},pdfauthor={Research prepared with AI assistance for Vladimir Reshetnikov}}
\newcommand{\N}{\mathbb N}
\newcommand{\Z}{\mathbb Z}
\newcommand{\Q}{\mathbb Q}
\newcommand{\Con}{\operatorname{Con}}
\newcommand{\ZFC}{\mathrm{ZFC}}
\newcommand{\ZF}{\mathrm{ZF}}
\newcommand{\code}[1]{\nolinkurl{#1}}
\newtheorem{theorem}{Theorem}[section]
\newtheorem{lemma}[theorem]{Lemma}
\newtheorem{proposition}[theorem]{Proposition}
\newtheorem{corollary}[theorem]{Corollary}
\theoremstyle{remark}
\newtheorem{remark}[theorem]{Remark}
\numberwithin{equation}{section}
\setlength{\parindent}{0pt}
\setlength{\parskip}{5pt}
\setlength{\emergencystretch}{2em}
\setlist{itemsep=2pt,topsep=3pt}
\pagestyle{fancy}
\fancyhf{}
\fancyhead[L]{\small Diophantine consistency: centered residues}
\fancyhead[R]{\small 25 October 2026}
\fancyfoot[C]{\thepage}
\setlength{\headheight}{14pt}
\lstset{basicstyle=\ttfamily\footnotesize,breaklines=true,columns=fullflexible,keepspaces=true,frame=single,rulecolor=\color{black!20}}
% Generated from arithmetic/signed258/metrics.json.
\newcommand{\SignedVariables}{258}
\newcommand{\SignedRows}{186}
\newcommand{\SignedTerms}{5,366}
\newcommand{\SignedHeight}{128,446}

\begin{document}
\hypersetup{pageanchor=false}
\begin{titlepage}
\vspace*{14mm}
{\Large Research contribution for ProveIt\par}
\vspace{9mm}
{\huge\bfseries Centered residue certificates\par}
\vspace{3mm}
{\huge\bfseries and reusable Pell denominators\par}
\vspace{8mm}
{\Large Smaller explicit Diophantine equations\par}
{\Large for the syntactic consistency of ZFC\par}
\vspace{12mm}
Prepared for Vladimir Reshetnikov\par
25 October 2026\par
Research, code, and exposition developed with AI assistance\par
\vspace{11mm}
\begin{abstract}
We give new explicit integer-coefficient polynomials whose absence of solutions, on their stated domains, is provably equivalent in ZFC to $\Con(\ZFC)$. The main structural reduction replaces eleven separate signed remainder intervals by quadratic radius certificates. A small canonical-residue bound, established before the factorial interpretation, makes these certificates sufficient to recover all the original Euclidean remainders. Two product coordinates merge the first two groups, and a parity construction replaces a three-square radix guard by one odd-square variable. The resulting direct proof-calculus quartic has \SignedVariables{} unrestricted integer variables, \SignedRows{} quadratic residuals, \SignedTerms{} collected monomials, and coefficient height \SignedHeight{}.

A separate admissible-denominator theorem improves the compact native construction. Its ten-positive-variable equations have operation/degree pairs $75/626$, $76/464$, and $77/460$. Every large coefficient is supplied literally and by a complete finite reconstruction. Natural-domain companion quartics and all coefficient expansions are included. The new proofs are ordinary mathematical proofs composed with the identified inherited numerical-calculus and native calibrations; the package does not assert a new end-to-end proof-assistant verification or a global simplicity minimum.
\end{abstract}
\vfill
\small
Comparison snapshot: ProveIt commit\par
\code{e0c98c629bf10d30b695884f6b1d2ab361bdf513}.\par
All six ZIP archives then present in \code{docs/incoming} were inspected.
\end{titlepage}
\pagenumbering{roman}
\hypersetup{pageanchor=true}
\begingroup
\small
\setlength{\parskip}{0pt}
\tableofcontents
\endgroup
\clearpage
\pagenumbering{arabic}

\section{The result and what ``simplest'' means here}

We use $\N=\{0,1,2,\ldots\}$. A consistency statement refers to the absence of a finite syntactic derivation of falsity from the repository's fixed ZFC axiom predicate. No semantic soundness assumption about ZFC is imposed. The aim is a fixed polynomial $P$ and a fixed witness domain $D$ such that
\begin{equation}\label{eq:target}
\ZFC\vdash \bigl(\neg\exists\mathbf x\in D\;P(\mathbf x)=0\bigr)
\longleftrightarrow \Con(\ZFC).
\end{equation}
An equation of this kind does not decide its own solvability. In particular, the proof of \eqref{eq:target} does not prove $\Con(\ZFC)$.

There is no single accepted ordering of ``simplest.'' Number of unknowns, degree, expanded support, coefficient height, binary evaluation cost, and the size of a complete coefficient specification measure different things. This contribution gives improvements on two concrete fronts, and records their domains separately.

% Generated comparison table; source identities are in PROVENANCE.json.
\begin{table}[htbp]
\centering
\small
\begin{tabular}{@{}llrrrr@{}}
\toprule
Construction & Domain & Variables & Degree & Rows & Monomials \\
\midrule
Incoming cofinal direct & $\Z$ & 282 & 4 & 192 & 5,491 \\
New signed principal & $\Z$ & 258 & 4 & 186 & 5,366 \\
New signed companion & $\Z$ & 259 & 4 & 185 & 5,361 \\
Incoming pooled natural & $\N$ & 226 & 4 & 175 & 10,314 \\
New natural balanced & $\N$ & 222 & 4 & 172 & 9,522 \\
New natural dimension & $\N$ & 221 & 4 & 171 & 10,216 \\
New natural sparse & $\N$ & 223 & 4 & 173 & 6,858 \\
\bottomrule
\end{tabular}
\caption{Direct numerical-calculus quartics. The three signed rows have
coefficient height 128,446. The natural heights and tradeoffs are given in
Section~\ref{sec:natural}. Comparators are the incoming packages
\cite{cofinal,graphdigits}, not a global frontier over all encodings.}
\end{table}

\begin{table}[htbp]
\centering
\small
\begin{tabular}{@{}lrrrr@{}}
\toprule
Compact construction & Positive variables & Operations & Degree & Denominator \\
\midrule
Incoming square transport & 10 & 78 & 630 & $Y^2$ \\
Incoming product transport & 10 & 79 & 468 & $Y^2$ \\
New square transport & 10 & 75 & 626 & $XY$ \\
New product transport & 10 & 76 & 464 & $XY$ \\
New lower-degree product & 10 & 77 & 460 & $q^4$ \\
\bottomrule
\end{tabular}
\caption{Exact binary $+,-,\times$ evaluation costs, with shared
intermediates counted once and squaring counted as a multiplication.
The coherent fixed coefficient tuple is unchanged from the incoming
compact parent \cite{certificate}.}
\end{table}


The strongest direct integer parent in the newly arrived cofinal-anchor package has 282 unknowns, 192 quadratic residuals, 5,491 monomials, and height 128,446 \cite{cofinal}. The new direct equation uses the same numerical proof-calculus route. It does not read or simulate K215 and does not specialize Jones's universal equation. The compact equations retain the established indexed native tuple; their improvements concern the arithmetic certificate, rather than a new logical encoding.

These are not global variable-count records. The audited repository already contains nine-positive-variable equations and separate universal-polynomial quartics with far fewer unrestricted integer variables than the direct construction \cite{repository}. The direct construction has a different attraction: all coefficients are small explicit integers, and its mathematical route starts from a finite numerical proof calculus. Our comparison of direct quartics is therefore explicitly restricted to that family.

\subsection{Precise finite definitions}

Each direct equation is defined by its complete ordered list of integer quadratic residuals:
\begin{equation}\label{eq:sos}
P_{\mathrm{dir}}(x_0,\ldots,x_{\SignedVariables-1})
=\sum_{j=0}^{\SignedRows-1} R_j(x_0,\ldots,x_{\SignedVariables-1})^2.
\end{equation}
Every $R_j$ is printed in the appendix and supplied in \code{arithmetic/signed258/equation.json}; the independently collected polynomial is \code{arithmetic/signed258/expanded_polynomial.json}. An entry $[c,[i_1,\ldots,i_d]]$ denotes $c x_{i_1}\cdots x_{i_d}$, with the empty list denoting the constant term. The JSON is a literal finite coefficient list, not an oracle, a universal-polynomial placeholder, or a variable-exponent expression.

The compact equation is printed in Section~\ref{sec:compact}. All of its auxiliary names are polynomial abbreviations. Precisely ten names are existentially quantified. Its four very large coefficients are fixed integers, supplied as byte payloads together with a finite source recipe. They are not additional unknowns.

\section{A centered-residue certificate}\label{sec:radius}

The main new idea is to exploit information that the power representation already supplies. Once a modular numerator has been interpreted, its canonical low residue is known to be very small relative to its modulus. Consequently it suffices to put an unrestricted representative in a short interval around zero. One quadratic bound can do this for several representatives simultaneously.

\begin{lemma}[Centered uniqueness]\label{lem:center}
Let $D\ge1$ be an integer. Suppose $r,r_0\in\Z$, $r\equiv r_0\pmod D$, $0\le r_0<D/2$, and $|r|<D/2$. Then $r=r_0$.
\end{lemma}
\begin{proof}
The difference $r-r_0$ is a multiple of $D$ and belongs to $(-D,D/2)$. Its only possible multiple of $D$ is zero.
\end{proof}

Put
\begin{equation}\label{eq:g}
g(a,b,c)=a^2+b^2+c(c+1).
\end{equation}
For integer inputs this is nonnegative, and every nonnegative integer is a value of $g$. For completeness of the latter statement, apply Legendre's three-square theorem to $4n+1$. It is not of the excluded form $4^e(8m+7)$. A three-square representation has exactly one odd coordinate; writing the even coordinates as $2a,2b$ and the odd one as $2c+1$ gives $n=g(a,b,c)$. This same standard representation is used in the inherited signed construction \cite{cofinal,sun}.

\begin{theorem}[One radius certificate for several residues]\label{thm:radius}
Let $1\le m\le56$, $D>0$, and $r_1,\ldots,r_m\in\Z$. Suppose the surrounding equations establish congruences
\[
r_i\equiv r_i^0\pmod D,\qquad 0\le r_i^0<D/15.
\]
The single quadratic equation
\begin{equation}\label{eq:radius}
D^2-4\sum_{i=1}^m r_i^2-1-g(a,b,c)=0
\end{equation}
forces $r_i=r_i^0$ for every $i$. Conversely the canonical representatives extend to a solution of \eqref{eq:radius} with integer $a,b,c$.
\end{theorem}
\begin{proof}
Since $g\ge0$, equation \eqref{eq:radius} gives $4r_i^2\le D^2-1<D^2$. Thus $|r_i|<D/2$, and Lemma~\ref{lem:center} identifies every representative. Conversely,
\[
4\sum_i(r_i^0)^2<\frac{4m}{225}D^2\le\frac{224}{225}D^2<D^2.
\]
The integer $D^2-4\sum_i(r_i^0)^2-1$ is therefore nonnegative and has a representation \eqref{eq:g}.
\end{proof}

The certificate has degree two. Squaring it as one residual in \eqref{eq:sos} preserves total degree four. In particular, there is no hidden appeal to a nonpolynomial norm, an integer square-root operation, or a higher-degree sign test.

\subsection{Why this replaces the old interval witnesses}

The parent supplied, separately for eleven remainders, a slack $s_i$ and a three-variable representation of a nonnegative product:
\[
r_i+s_i=D_i-1,\qquad r_i s_i=g(a_i,b_i,c_i).
\]
Once $D_i>0$ was known, these equations implied $r_i,s_i\ge0$, hence $0\le r_i<D_i$. The eleven moduli belong to four groups of sizes $5,2,2,2$. Equation \eqref{eq:radius} replaces each whole group. Given its solution, set $s_i=D_i-1-r_i$ and restore the old product witnesses. Conversely, the arithmetic bounds below show that the parent already used the small canonical residues, so a new radius gap is representable.

This is an equivalence of existential zero relations. It does not identify the old product witnesses with the new radius witnesses. The literal transformation deletes only coordinates private to the removed constraints, retains every unrelated row, and supplies explicit maps for the remaining coordinates. The final projection is checked separately; the number of deleted source witnesses alone does not predict the size of the projected quartic.

\subsection{Merging two adjacent-power groups}

The two low remainders at modulus $D_t=q^t$ and the five at modulus $D_{t+1}=q^{t+1}$ can share one radius. Introduce just two product coordinates
\[
u_1=qr_1,\qquad u_2=qr_2.
\]
Use $u_1,u_2$ and the other five remainders in \eqref{eq:radius}, with modulus $D_{t+1}$. Power recovery already gives $D_{t+1}=qD_t$ before this radius is interpreted. The canonical scaled remainders obey $q r_i^0<D_{t+1}/15$. Centered uniqueness therefore gives $u_i=qr_i^0$, and division by the positive integer $q$ gives $r_i=r_i^0$.

The merged group has seven members, well within Theorem~\ref{thm:radius}. The two product equations are quadratic, so the final polynomial is still quartic. This change saves one variable relative to two separate three-variable radius gaps, while adding one residual. The final construction has groups of sizes $7,2,2$; the separate-group companion remains available because it has fewer rows and slightly fewer collected monomials.

\section{The small-residue bound before factorial semantics}\label{sec:bootstrap}

Here we state the exact inherited arithmetic interface on which the new reduction acts. Full predecessor proofs are retained beside the delivered source. The signed frame and orientation argument establishes the canonical Pell polynomials at each requested signed index before interpreting the power requests. Its square frame includes the actual bases, indices, and outputs, so this preliminary step uses no positivity assumption on the remainder variables being removed \cite{cofinal}.

\subsection{Power and endpoint recovery}

Write $t=16+u^2$, and let $F,B,V_0,x,y$ denote the parent's semantic coordinates. Its two retained factorial remainder guards give $B\ge1$ and
\begin{equation}\label{eq:ceil}
FB=V_0+s,\qquad s\ge0.
\end{equation}
The correctly oriented request at the positive index $t$ yields
\[
V_0=(F^2+1)^t\ge1.
\]
Thus $F\ge1$, $B>F$, and $B\ge2$. The cofinal beta base is
\[
y=V_0(V_0-1)>F+2.
\]
The signed power representation at base $B\ge2$ implies that its requested exponent $x$ is nonnegative. Indeed a negative exponent would give, by the same large-modulus congruence bound as the parent, an impossible equality $mB^j=1$ with $j\ge1$ and integer $m$. The two endpoint equations then imply
\begin{align}
x&=F+1+(y+1)u_0,\\
x&=1+[1+(t+1)y]u_t,\\
u_0&\ge0,\quad u_t\ge1,\quad x\ge2+(t+1)y.\label{eq:xlower}
\end{align}

Let $a_1,\ldots,a_9$ be the actual independent CRT representatives, and set
\[
H_0=1+\sum_{i=1}^9 a_i,\qquad K=x,\qquad U=BH_0.
\]
The radix guard supplies $q\ge1$. Each representative is the exponent of a request at base $q+1\ge2$, so the same negative-index argument forces $a_i\ge0$. It follows that $H_0\ge1$. The retained requests now establish
\begin{equation}\label{eq:powers}
N=B^K,\qquad q=(BH_0)^K=U^K,
\qquad q^t,\ q^{t+1},\ q^{y+1},\ q^x.
\end{equation}
The derived-label numerator has upper index $J\ge0$, obtained from the fixed numerical constructor polynomial in the four parameter representatives. Its zero case is already excluded by the semantic Pell anchor. All numerator requests are consequently interpreted before any low-digit interval is used.

For the radius-only intermediate, the old guard supplies $q\ge2$. For the final odd-square variant, Section~\ref{sec:odd-square} supplies $q\ge1$ and the recovery order just given. Both yield $q\ge B^2>V_0$ after \eqref{eq:powers}.

\subsection{An explicit margin}

Let $\ell=2^{22}$ be the unchanged checked degree majorant of the label polynomial. The source's constructor bound gives
\begin{equation}\label{eq:jbound}
J\le(\max(1,a_1,a_2,a_3,a_4)+1)^\ell\le(2H_0)^\ell\le U^\ell.
\end{equation}
This estimate is valid even when some actual representative is zero; it does not assume the stronger post-factorial positivity of the encoded index.

Before identifying $F$ as a factorial, $V_0\ge2^t$ and \eqref{eq:xlower} give
\[
K\ge2+(t+1)2^t(2^t-1)>\ell(t+1)+4\qquad(t\ge16).
\]
For example, $2^{16}(2^{16}-1)>2^{22}+4$; the left side increases with $t$, so this elementary bound persists. In particular $K\ge6$ and $q\ge2^K>16$.

For an ordinary upper index $a_i\le H_0$, every coefficient needed strictly below a lower index $k\le K$ satisfies
\[
\binom{a_i}{j}\le\max(1,a_i)^j\le H_0^{K-1}
\le\frac{q}{B^K}\le\frac q{16}\qquad(0\le j<k).
\]
For the derived-label request, $k=t+1$ and \eqref{eq:jbound} gives
\[
\binom Jj\le\max(1,J)^j\le U^{\ell t}
\le U^{K-4}\le q/16.
\]
For the factorial/product request, the inherited elementary binomial bound is
\[
\binom{F^2+t}{j}\le(F^2+1)^t=V_0<B^2\le q/16\qquad(0\le j\le t).
\]
The last inequality uses $q\ge B^K$, $B\ge2$, and $K\ge6$. This is deliberately established without invoking $F=t!$.

\begin{lemma}[Low digits have room]\label{lem:digits}
Suppose $q\ge16$, $D=q^k$, and all binomial coefficients through degree $k-1$ are at most $q/16$. The canonical remainder of $(q+1)^n$ modulo $D$ is
\[
r^0=\sum_{j=0}^{k-1}\binom nj q^j,
\]
and $0\le r^0<D/15$.
\end{lemma}
\begin{proof}
Every displayed coefficient is below $q$, so the expression is a valid base-$q$ block. Moreover
\[
r^0\le\frac q{16}\frac{q^k-1}{q-1}<\frac{D}{15},
\]
using $q/(16(q-1))\le1/15$ for $q\ge16$ and the strict inequality $D-1<D$.
\end{proof}

All eleven low remainders therefore meet Theorem~\ref{thm:radius}'s hypotheses before factorial semantics is used. The four separate radius equations, or the three equations after the weighted merge, recover their exact canonical values. The modular digit for $\binom{F^2+t}{t}$ now gives the exact output $B$; the unchanged ceiling fixed-point argument gives $F=t!$. The original CRT range extraction and Horn-flow proof then apply in their established order. This ordering is the substantive point: a later interpretation is never used to justify a bound needed to obtain that interpretation.

\section{One odd square can supply the initial radix guard}\label{sec:odd-square}

The radius-only intermediate still retains the old exact-binomial output slack $o=q-B-1$ together with its three-variable sign certificate. The final construction removes that private slack and its sign witnesses and imposes instead
\begin{equation}\label{eq:oddq}
q=(2z+1)^2,
\end{equation}
where $z$ is one unrestricted integer. This is a quadratic residual and guarantees $q\ge1$ without another sign guard.

\subsection{Soundness}

Use the order in Section~\ref{sec:bootstrap}: recover $V_0$ at the already positive index $t$, then $F\ge1$ and $B\ge2$, then $x\ge0$ from the request $N=B^x$, and finally \eqref{eq:xlower}. The base $q+1\ge2$ orients all representative indices. Consequently $H_0\ge1$ and the exact radix request yields
\[
q=(BH_0)^x\ge B^2>B.
\]
Thus the omitted slack $o=q-B-1$ is nonnegative and admits a fresh representation \eqref{eq:g}. The remaining parent proof can be restored. In particular, neither $q>1$ nor $B<q$ was silently assumed to prove the exact power that supplies it.

\subsection{Completeness and simultaneous parity choices}

Start with a finite refutation and the parent's cofinal horizon and beta construction. Its chosen $y=V_0(V_0-1)$ is even. Every beta modulus $1+(i+1)y$ is therefore odd; their common period is odd. Adding that period permits choosing the global beta code $x$ even, while preserving all local residues and making it as large as required.

At $F=t!$ the reciprocal CRT factors are $d_i=1+F^2/i$. Since $i\mid F$ and $F$ is even,
\[
F^2/i=F(F/i)\quad\hbox{is even},
\]
so every $d_i$, their product $B$, and $N=B^x$ are odd. CRT representatives can therefore be chosen with arbitrary prescribed parity by shifting them by the odd period $N^2$.

Keep the parent's opposite parity for the two representatives determining the semantic anchor's parity. Choose the parity of a different representative to make $H_0=1+\sum a_i$ odd. This second choice does not change the anchor's parity. Ordinary representatives remain in $(0,2N^2]$; the designated dominant representative stays in $[2N^2,4N^2)$ and retains the parent's dominance bound. Thus all required anchor and range conditions still have witnesses.

Now $U=BH_0$ is odd and $x$ is even. Hence
\[
q=U^x=(U^{x/2})^2,
\qquad z=\frac{U^{x/2}-1}{2}\in\Z
\]
satisfies \eqref{eq:oddq}. The radius gaps are nonnegative by Theorem~\ref{thm:radius}, so fresh three-square witnesses finish the construction. All choices occur in the original dependency order; in particular the self-sized radix is chosen after the CRT representatives.

\section{The direct polynomial and the logical endpoint}\label{sec:logical}

The raw signed construction is transformed by exact algebraic factorization, elimination of integer coordinates with monic defining equations, and an integer translation. Each retained elimination has a polynomial restoration map. The independent checker verifies all substitutions and the complete collected sum of squares. The translation preserves the integer zero set and reduces coefficient size. These finite operations account for the final measured \SignedVariables{} coordinates, rather than an informal subtraction from the parent's count.

The remaining arithmetic source is the inherited numerical calculus $\mathcal C_{222}$. Its finite local Horn-flow conditions recover a derivation of contradiction label $1$. Conversely any finite derivation supplies a padded local flow, CRT representatives, power witnesses, and the present radius and parity witnesses. Write $\operatorname{Ref}_{T}$ for the existence of a finite syntactic refutation in theory $T$. The two logical interfaces retained in the package are
\begin{align*}
\mathcal C_{222}(1)&\longleftrightarrow\operatorname{Ref}_{\ZF},\\
\operatorname{Ref}_{\ZF}&\longleftrightarrow\operatorname{Ref}_{\ZFC}.
\end{align*}
They are the repository theorems \code{KEndpoint.c222_iff_zf_refutation} and \code{ConstructibleChoice.zf_refutation_iff_zfc_refutation} \cite{endpoint,choice}. The first is a machine-independent statement even though its file also contains machine endpoints. The second uses the constructible-model argument for removal of Choice.

\begin{theorem}[The new direct equation]\label{thm:main}
Let $P_{\mathrm{dir}}$ be the fixed polynomial \eqref{eq:sos} with the literal coefficients supplied in this package. Its degree is four, its domain is $\Z^{\SignedVariables}$, and
\[
\ZFC\vdash
\left[\neg\exists\mathbf x\in\Z^{\SignedVariables}\;
P_{\mathrm{dir}}(\mathbf x)=0\right]
\longleftrightarrow\Con(\ZFC).
\]
\end{theorem}
\begin{proof}
Over the integers a sum of squares is zero exactly when every residual is zero. The restoration maps return a zero of the changed arithmetic source. Sections~\ref{sec:radius}--\ref{sec:odd-square} prove its equivalence to the calibrated finite-refutation relation, in both directions. Compose this equivalence with the two logical interfaces and negate it. All finite sequences, integer inequalities, Pell inductions, CRT constructions, and finite syntax used here are formalizable in ZFC.
\end{proof}

The statement is an ordinary formalizable theorem, rather than a newly emitted derivation in the object language of ZFC. The source interfaces are retained for inspection; this package does not rebuild their transitive Lean dependency graph. The new algebra checkers verify the emitted finite objects and restoration maps, not the truth of every unbounded mathematical claim by testing.

\section{Natural-domain companion quartics}\label{sec:natural}

There is a separate useful composition on the nonnegative-integer route.
It combines the cofinal semantic anchor with the four-group low-digit
budget of the graph-digits package \cite{cofinal,graphdigits}.  Its three
selected endpoints are
\begin{center}
\begin{tabular}{@{}lrrrr@{}}
\toprule
Variant & Variables & Residuals & Monomials & Height\\
\midrule
Smaller dimension & 221 & 171 & 10,216 & 2,445,412,949,744\\
Balanced & 222 & 172 & 9,522 & 83,954,186\\
Sparse & 223 & 173 & 6,858 & 83,953,724\\
\bottomrule
\end{tabular}
\end{center}
All have exact degree four and range over $\N$, not over unrestricted
integers. The balanced and sparse points avoid the very large coefficient
height paid by the 221-variable projection. The balanced 222-variable
point improves every listed metric of the incoming 226-variable
graph-digits endpoint, which had 175 residuals, 10,314 monomials, and
height 87,839,583.

\subsection{Pooling bounds in a nonnegative domain}

Suppose $m$ nonnegative remainders share a positive denominator $D$.
Instead of $m$ separate rows $r_i+s_i=D-1$, retain one nonnegative slack
$s$ and impose
\begin{equation}\label{eq:nat-pool}
\sum_{i=1}^m r_i+s=D-1.
\end{equation}
Every individual interval is recovered by
$s_i=s+\sum_{j\ne i}r_j\ge0$. Completeness requires the additional
bound $\sum_i r_i<D$. This is a genuine extra requirement; it cannot be
assumed merely because each $r_i<D$ separately.

The natural parent's radix is
\[
q=(NS)^K,\qquad K=x,\qquad
S=1+C+\sum\text{the eight independent encoded columns}.
\]
This $S$ is distinct from the signed construction's $H_0$, and the base
contains $N$, not $B$. Let $U=NS$. Every ordinary positive upper index is
less than $S$, so each of its needed low coefficients is at most
$U^{K-1}$. The derived label satisfies $J\le S^\ell$ and
$\ell t<K$, giving the same bound for its coefficients through degree
$t$. In canonical completeness the factorial upper index $F^2+t<N$;
its low coefficients also obey the same bound. The inherited inequalities
give $U>5$.

For a common lower index $k$, the sum of the $m\le5$ canonical low
blocks therefore obeys
\[
\sum_i r_i\le mU^{K-1}\frac{q^k-1}{q-1}<q^k=D.
\]
For the strict inequality, $m\le5<U$ implies
$mU^{K-1}\le q-U^{K-1}\le q-1$. Thus \eqref{eq:nat-pool} has its
required nonnegative slack. Four groups replace eleven individual
bounds, saving seven source slack coordinates and seven rows.

There is also a redundant output bound. Immediately after exact power
recovery, the ceiling gives $F\ge1$, $B>F$, and $B\ge2$; the endpoints
give $K=x\ge2$. Hence
\[
N=B^K\ge B^2,\qquad q=(NS)^K\ge N^2\ge B^4>B.
\]
The output slack $q-B-1$ can be restored nonnegatively without using the
factorial conclusion. Deleting it changes the later projection and
slightly improves coefficients, but is dimension-neutral after the
selected contractions. The final numbers in the table are measured
from the complete emitted polynomials.

\subsection{Proof and exact projection}

A new natural zero restores every old remainder slack by the displayed
nonnegative polynomial formula and restores the output slack using the
power bootstrap. The inherited cofinal-anchor source then decodes a
finite refutation. Conversely, the canonical refutation witnesses have
the spare radix capacity proved above and extend to the pooled rows.
The different final projections have checked natural restoration maps;
two of the restorations use exact rational sum-of-squares certificates
for nonnegativity rather than an incorrect coefficientwise sign test.

Each natural companion therefore satisfies the same ZFC-provable
no-zero equivalence on its specified natural domain. The complete
balanced residual list appears in Appendix~\ref{app:natural}; all three
literal expansions and their certificates are included in the archive.

\section{Compact equations from an existing denominator}
\label{sec:compact}

The second contribution is a family of explicit polynomials in ten
strictly positive integer unknowns.  It improves the latest incoming
\(78/630\) and \(79/468\) circuits, where the two numbers mean live binary
arithmetic operations and exact total degree.  The new arithmetic permits
an existing expression to serve as the carrier denominator; two further
savings reuse products already needed elsewhere.  The complete results
are as follows.

\begin{center}
\begin{tabular}{@{}llrr@{}}
\toprule
Denominator & Transport weight & Operations & Exact degree\\
\midrule
\(XY\) & \(V^2\) & \textbf{75} & \textbf{626}\\
\(XY\) & \(bV\) & \textbf{76} & \textbf{464}\\
\(q^4\) & \(V^2\) & 76 & 622\\
\(q^4\) & \(bV\) & \textbf{77} & \textbf{460}\\
\bottomrule
\end{tabular}
\end{center}

The \(76/622\) point has the simpler squared-weight proof, although
\(76/464\) dominates it in the two numerical metrics.  The other three
points trade arithmetic cost against degree.  All four use the same
complete indexed BFS943 coefficient tuple, specified in
Section~\ref{sec:compact-constants}; its four ports are constants, not
additional unknowns.  The arithmetic changes are uniform in those ports
within their stated box, but the consistency interpretation uses the
calibration of that \emph{whole} tuple.\cite{certificate,indexed}

\begin{theorem}[Compact consistency equations]
\label{thm:compact-main}
Each of the four fixed polynomials defined below has ten strictly positive
integer unknowns.  Composing the ordinary arithmetic proof in this
section with the retained indexed native/refutation calibration gives
\[
 \ZFC\ \vdash\
 \bigl[\nexists\,\mathbf x\in\Z_{>0}^{10}:P(\mathbf x)=0\bigr]
 \ \longleftrightarrow\ \Con(\ZFC).
\]
Their operation counts and exact total degrees are those in the table.
\end{theorem}

This is an ordinary theorem formalizable in ZFC.  It is not a new
end-to-end Lean or Rocq derivation, an emitted object-language ZFC proof,
a proof of consistency, or a claim of global optimality.  Pell doubling,
the product-square finalizer, and the removal of the auxiliary shift have
explicit predecessors.  The new denominator criterion and its completed
composition are distinguished from those inherited steps below.
\cite{certificate,unshift}

\subsection{The complete polynomial}
\label{subsec:compact-definition}

The independent unknowns are exactly
\[
 J,F,\alpha,i,T,w,s,n,\eta,\zeta\in\Z_{>0}.
\]
The four fixed positive integers are \(\beta,K_0,M_C,M_F\), with
\begin{equation}
 \beta\ge31\text{ odd},\qquad K_0>0,\qquad
 0<M_C<\beta<M_F<2\beta.
 \label{eq:compact-portbox}
\end{equation}
Define the following polynomial abbreviations:
\begin{align}
 h&=\beta J,& q&=h+1,& C&=q-F-\alpha,\nonumber\\
 X&=wq,&Y&=q^3s,&k&=\eta+\zeta,\nonumber\\
 c&=kY+\eta,& A&=Y(X+1)+1,&
 \Delta&=A^2-1,\nonumber\\
 H&=4A-5,& \ell&=(XY)(kY),&
 E&=1+\ell(\ell+k).                         \label{eq:compact-base}
\end{align}
For the mask, transport, and strong norm, put
\begin{align}
 R&=[h(q-F)+\alpha](q^2-1)+(M_C+qM_F)J,\nonumber\\
 e&=k-R-1, & B&=(K_0+w)C-F,\nonumber\\
 U&=ic^2,& a&=\Delta U^2E,& F_0&=1+a,\nonumber\\
 G&=EF_0,& \mathcal F&=1+4aF_0,&
 Q&=\Delta(\mathcal F-1),\nonumber\\
 b&=cT-R,& V&=\mathcal F b-c,&
 K_+&=X(3c-2X)+2.                            \label{eq:compact-strong}
\end{align}
Here \(\mathcal F=(1+2a)^2\) is an identity; its square root has not
been introduced as an unknown.  Choose one of
\[
 D=XY\quad\hbox{or}\quad D=q^4,
 \qquad
 L=V^2\quad\hbox{or}\quad L=bV.
\]
Finally set
\begin{align}
 M&=hDH,\nonumber\\
 Z&=hD\bigl[H(n-LB)+K_+\bigr]-H(DB+he),\nonumber\\
 R_G&=M^2G.                                             \label{eq:compact-carrier}
\end{align}
The requested single equation is
\begin{equation}
 \boxed{\ P=Q(Z^2-R_GV^2)+(R_G-Z^2)=0.\ }
 \label{eq:compact-final}
\end{equation}
Every named quantity in these displays is a polynomial abbreviation.
The delivered JSON graphs contain only binary addition, subtraction, and
multiplication.  Squaring and multiplication by fixed scalars are counted
as multiplications.  No division or variable exponent occurs in the
polynomial.

The displayed form also makes the two shared products visible:
\(A-1=XY+Y\), while \(kY\), already used in \(c\), is reused in
\(\ell=(XY)(kY)\).  The proof occasionally writes \(V_0=XY^2\), but
the evaluator need not compute that expression separately.

\subsection{The admissible-denominator criterion}
\label{subsec:compact-denominator}

The same proof works with a more general integer-polynomial denominator
\(D\), provided that, at every positive input,
\begin{equation}
 D>0,\qquad q^4\mid D,\qquad D\mid4XY^2.
 \label{eq:compact-admissible}
\end{equation}
The first divisibility will recover a bounded index error.  The second
will recover the complete quotient only \emph{after} that error is known
to vanish.  Pairwise coprimality of \(D\) with all the other denominators
is unnecessary.  Both advertised choices satisfy the criterion:
\begin{align*}
 XY&=q^4(ws),&4XY^2&=(XY)(4Y),\\
 4XY^2&=q^4(4ws^2q^3).
\end{align*}
In particular, \(D\ge q^4>2\).

Both choices depend only on the outer coordinates \(J,w,s\), so the
canonical completeness construction below has no circular choice of a
padding coordinate.  In fact the pointwise criterion already excludes
such a dependence for any polynomial \(D\): with \(J,w,s\) fixed, it
implies \(0<D\le4XY^2\) while any of the other coordinates varies through
all positive integers.  A bounded polynomial in such a coordinate is
constant.  Repeating this argument removes every non-outer coordinate.
Within this criterion, total degree four is minimal for a polynomial
denominator: fix \(w,s\) and let \(J\) grow, using \(D\ge(\beta J+1)^4\).
The choice \(q^4\) attains that denominator-degree bound; this is not a
global bound on consistency equations.

The rational quantities used only in the proof are
\[
 z=\frac ZM,\qquad \kappa=\frac{K_+}{H},\qquad
 \rho=\frac eD,\qquad t=\frac Bh,\qquad W=1+hL.
\]
The exact identities are
\begin{align}
 z&=n+\kappa-Wt-\rho,                         \label{eq:compact-rational}\\
 P&=(Q-1)Z^2-M^2G(QV^2-1).                   \label{eq:compact-norm}
\end{align}
No divisibility premise is used in either identity.  In particular,
\(\rho\) and \(t\) are not initially assumed integral.

\subsection{Pell facts and the index lemmas}
\label{subsec:compact-pell}

For an integer \(A\ge2\), define
\[
 \chi_A(r)+\psi_A(r)\sqrt{A^2-1}
       =(A+\sqrt{A^2-1})^r\qquad(r\ge0).
\]
Both sequences satisfy \(u_{r+2}=2Au_{r+1}-u_r\), with
\(\chi_A(0)=1,\chi_A(1)=A\) and
\(\psi_A(0)=0,\psi_A(1)=1\).
Positive integer solutions of \(u^2-(A^2-1)v^2=1\) are exactly these
pairs.  One elementary descent proves this classification: for \(v>1\),
multiplication by \(A-\sqrt{A^2-1}\) gives a positive integer ordinate
\(Av-u\) strictly smaller than \(v\); at \(v=1\), \(u=A\).
The addition formulas also give strong divisibility,
\begin{equation}
 \gcd\bigl(\psi_A(r),\psi_A(t)\bigr)
       =\psi_A(\gcd(r,t)).                         \label{eq:compact-gcd}
\end{equation}
Indeed, \(\chi_A(r)\) is a unit modulo \(\psi_A(r)\), and the addition
law reduces the gcd of two ordinates exactly as the Euclidean algorithm
reduces their indices.  These elementary facts are the Pell framework
of the retained normalized kernel.\cite{native61}

We also need the first norm in \eqref{eq:compact-base}.  If \(v\ge1\)
and positive integers \(\tau,k\) satisfy
\(\tau^2-v(v+1)k^2=1\), then
\begin{equation}
 \tau=\chi_{2v+1}(r),\qquad k=2\psi_{2v+1}(r)
 \quad\hbox{for some }r\ge1.                       \label{eq:compact-first-class}
\end{equation}
For completeness, there is no solution with \(k=1\), because
\(v(v+1)+1\) lies strictly between consecutive squares.  The pair
\((2v+1,2)\) is a solution.  For \(k>2\), multiplication by its inverse
unit gives the positive smaller ordinate \((2v+1)k-2\tau\).  Descent
therefore classifies every positive solution as a power of this pair,
which is precisely \eqref{eq:compact-first-class}.

\begin{lemma}[Odd-index quotient identities]
\label{lem:compact-odd}
There are integer polynomials \(Q_r\) such that
\begin{align*}
 \chi_t(2r+1)&=tQ_r(t^2),&
 Q_r(0)&=(-1)^r(2r+1),\\
 Q_r(1-A^2)&=(-1)^r\psi_A(2r+1),&
 \psi_t(2r+1)&\equiv(-1)^r\pmod{t^2}.
\end{align*}
\end{lemma}
\begin{proof}
Use \(Q_0(x)=1\), \(Q_1(x)=4x-3\), and
\(Q_{r+1}(x)=(4x-2)Q_r(x)-Q_{r-1}(x)\).
The odd trace subsequence gives the first identity.  Evaluation at zero
gives the second.  At \(x=1-A^2\), the recurrence coefficient is the
negative of the one for the odd \(\psi_A\) subsequence, whose initial
values are \(1,4A^2-1\); this gives the third identity.  The recurrence
for \(\psi_t(2r+1)\), with initial values \(1,4t^2-1\), gives the last
congruence modulo \(t^2\).
\end{proof}

\begin{lemma}[Normalized rank]
\label{lem:compact-rank}
Suppose \(c=\psi_A(p)>2\), \(f,I>0\),
\(f^2-(A^2-1)I^2=1\), and \(c^2\mid I\).
Then \(f=\chi_A(m)\), \(I=\psi_A(m)\) for an index satisfying
\(pc\mid m\), and hence \(c\mid m\) and \(m>2p\).
\end{lemma}
\begin{proof}
Classification gives \(m\), and \eqref{eq:compact-gcd} gives \(p\mid m\).
Write \(m=pv\) and \(D_0=\chi_A(p)\).  Composition gives
\(I=c\psi_{D_0}(v)\).  Since \(D_0^2\equiv1\pmod c\), recurrence
gives \(\psi_{D_0}(v)\equiv vD_0^{v-1}\pmod c\).
The number \(D_0\) is invertible modulo \(c\), and \(c^2\mid I\), so
\(c\mid v\).  Thus \(pc\mid m\) and \(m\ge pc>2p\).
\end{proof}

\begin{lemma}[Signed index recovery]
\label{lem:compact-index}
Let \(c=\psi_A(p)\), \(f=\chi_A(m)\), \(c\mid m\), \(m>2p\),
\(S\ge2\), \(c\mid S\), and \(S^2=(A^2-1)(f^2-1)\).
If \(0<p,R<c/2\) and integers \(V,y\) satisfy
\[
 S^2V^2-(S^2-1)y^2=1,\qquad
 V\equiv-c\pmod f,\qquad V\equiv-R\pmod c,
\]
then \(p=R\).  The signs of \(V,y\) are unrestricted.
\end{lemma}
\begin{proof}
Classification at parameter \(S\) gives
\(S|V|=\chi_S(\nu)\), \(|y|=\psi_S(\nu)\).
The index \(\nu\) is odd: an even trace is \(\pm1\pmod S\).
Lemma~\ref{lem:compact-odd}, using
\(S^2\equiv1-A^2\pmod f\), gives
\(\psi_A(\nu)\equiv\pm c\pmod f\).  Squaring and using
\(\chi_A(2j)=1+2(A^2-1)\psi_A(j)^2\) therefore gives
\[
 \chi_A(2\nu)\equiv\chi_A(2p)\pmod{\chi_A(m)}.
\]
Here is the abscissa step used to read this congruence.  For
\(0<k<m\), a congruence
\(\chi_A(t)\equiv\chi_A(k)\pmod{\chi_A(m)}\) implies
\(t\equiv\pm k\pmod{2m}\).  The addition laws reduce and reflect
\(t\) to a signed representative \(\pm\chi_A(u)\), \(0\le u\le m\).
The value \(u=m\) has residue zero.  The negative sign is impossible
because \(0<\chi_A(u)+\chi_A(k)\le2\chi_A(m-1)<\chi_A(m)\).
The positive sign forces equality and hence \(u=k\).
Apply this step with \(k=2p<m\): it gives
\(\nu\equiv\pm p\pmod m\), hence modulo \(c\).

On the other hand, \(c\mid S\) and \(Q_r(0)=(-1)^r(2r+1)\), together
with \(V\equiv-R\pmod c\), give \(\nu\equiv\pm R\pmod c\).
Thus \(p\equiv\pm R\pmod c\).  The strict half-interval bounds on
\(p,R\) force \(p=R\).
\end{proof}

\subsection{Root extraction before any transport sign}
\label{subsec:compact-roots}

At every positive input,
\[
 q\ge32,\quad X\ge q,\quad Y\ge q^3,\quad k\ge2,\quad c\ge3,
 \quad A>4,\quad\Delta\ge3,\quad U\ge c^2,\quad E>1.
\]
Moreover \(\mathcal F>2c\), \(Q>1\), and \(M>0\).
For every integral \(b\), \(V=\mathcal Fb-c\) is nonzero and
\(|V|>1\): if \(b=0\), \(V=-c\); otherwise
\(|V|\ge\mathcal F-c>1\).
The identities
\begin{equation}
 \mathcal F=(1+2a)^2,\qquad
 Q=4\Delta^2U^2G                                      \label{eq:compact-double}
\end{equation}
imply \(\gcd(G,Q-1)=\gcd(G,QV^2-1)=1\).
At a zero of \(P\), equation~\eqref{eq:compact-norm} yields
\[
 (Q-1)z^2=G(QV^2-1).
\]
For every prime dividing \(G\), its valuations in \(Q-1\) and
\(QV^2-1\) vanish; its valuation in \(G\) must therefore be even,
because the valuation of a rational square is even.  Thus \(G=g^2\)
for an integer \(g>0\).  Since \(F_0\equiv1\pmod E\), the coprime
positive factors \(E,F_0\) are separately squares:
\[
 E=\tau^2,\qquad F_0=f_0^2,\qquad g=\tau f_0.
\]
No coprimality condition on \(M\) was required.
Put \(y=|z|/g\).  Then
\begin{equation}
 QV^2-(Q-1)y^2=1.                                  \label{eq:compact-adjacent}
\end{equation}

\begin{lemma}[Integral ordinate in the adjacent norm]
\label{lem:compact-adjacent}
For an integer \(Q>1\), a nonzero integer \(V\), and a positive rational
\(y\), equation~\eqref{eq:compact-adjacent} forces \(y\) integral.
\end{lemma}
\begin{proof}
Write \(y=r/m\) in lowest terms.  Clearing denominators gives
\(m^2\mid Q-1\).  For \(v=|V|>1\), the norm implies
\(v<y<Qv/(Q-1)\).  The transformation
\[
 v'=(2Q-1)v-2(Q-1)y,\qquad
 y'=(2Q-1)y-2Qv
\]
preserves the norm.  Its first coordinate is integral and satisfies
\(-v<v'<v\); it cannot vanish, since the norm is one.
The reduced denominator of the second coordinate remains \(m\), because
\(2Q-1\equiv1\pmod m\).  Reflecting signs and descending on the positive
integer first coordinate eventually reaches \(v=1\), where the norm
gives \(y=1\).  Hence the unchanged denominator is \(m=1\).
\end{proof}

Consequently \(z=\varepsilon gy\) is an integer, with
\(\varepsilon\in\{-1,1\}\).  For proof purposes set
\begin{equation}
 f=1+2a,\qquad I=2Ug.
 \label{eq:compact-proofpair}
\end{equation}
Then
\begin{equation}
 f^2=\mathcal F,\qquad f^2-\Delta I^2=1,\qquad
 Q=(\Delta I)^2,\qquad c^2\mid I.                 \label{eq:compact-strongnorm}
\end{equation}

We next give the sign-free main-index folding argument explicitly.
Classify \(f=\chi_A(m)\), \(I=\psi_A(m)\).
At \(S=\Delta I\), classify \eqref{eq:compact-adjacent} as
\(S|V|=\chi_S(j)\), \(y=\psi_S(j)\).  Again \(j\) is odd.
Since \(S^2\equiv1-A^2\pmod f\),
Lemma~\ref{lem:compact-odd} and \(V\equiv-c\pmod f\) give
\(\psi_A(j)\equiv\pm c\pmod f\).  The addition laws give
\[
 \psi_A(j+2m)\equiv-\psi_A(j),\qquad
 \psi_A(2m-j)\equiv\psi_A(j)\pmod f.
\]
Reduction modulo \(2m\), followed by reflection about \(m\), supplies
\(p\in[0,m]\) of the parity of \(j\), with
\(\psi_A(j)\equiv\pm\psi_A(p)\pmod f\).
Now \(0<c\le I\), \(0\le\psi_A(p)\le I\), and \(2I<f\) because
\(A\ge3\).  The congruence therefore forces
\begin{equation}
 c=\psi_A(p),\qquad p\ge1\text{ odd}.              \label{eq:compact-folded}
\end{equation}
Neither the sign of \(V\) nor a parity condition on \(m\) was imposed.

Write \(V_0=XY^2\) for proof notation.  The first norm and
\eqref{eq:compact-first-class} give
\[
 k=2\psi_{2V_0+1}(n_0),\qquad n_0\ge1.
\]
If \(n_0=1\), then \(k=2\), \(\eta=\zeta=1\), and \(c=2Y+1\)
lies strictly between \(\psi_A(1)=1\) and \(\psi_A(2)=2A\), a
contradiction.  Thus
\begin{equation}
 n_0\ge2,\quad k\ge8V_0+4,\quad
 c>kY\ge k>8q^7,\quad c>X,\quad K_+>0.           \label{eq:compact-earlybound}
\end{equation}
No sign of \(B,C,R,b,V\), or \(\rho\) was used.

\subsection{Signed transport with no shift in \texorpdfstring{$T$}{T}}
\label{subsec:compact-sign}

The old coordinate was \(T+\alpha\).  Its formal inverse after deletion
is \(T_{\rm old}=T-\alpha\), which can be nonpositive.  Soundness therefore
needs a direct argument; a formal substitution alone would be inadequate.
\cite{unshift}

First \(b\ne0\).  Otherwise \(V=-c\), and
\eqref{eq:compact-adjacent} would give
\[
 y^2-c^2=\frac{c^2-1}{Q-1}\in(0,1),
\]
because \(Q\ge9c^4>c^2\).  The left side is integral.
Put \(D_0=\Delta U^2\ge243\).  Since \(a=D_0E\) and
\(g^2=E(1+a)\),
\[
 \frac{a^2}{g^2}=\frac{D_0a}{1+a}>\frac{D_0}{2}>9.
\]
Hence \(a>3g\), while \(a\ge c\).  For nonzero integral \(b\),
\begin{equation}
 \operatorname{sgn}V=\operatorname{sgn}b,\qquad
 |V|\ge\mathcal F-c>a>3g,\qquad |V|>|b|,\qquad
 y<\tfrac32|V|.                                    \label{eq:compact-absolute}
\end{equation}
For \(L=V^2\), these immediately give \(gy<L/2\).

The choice \(L=bV\) requires a stronger estimate.  Suppose the odd
main index in \eqref{eq:compact-folded} were \(p=3\), so
\(c=4A^2-1\).  If \(n_0\ge3\), then
\[
 k>32V_0^2,\qquad
 c<64X^2Y^2<32X^2Y^5<kY<c,
\]
using \(A<4XY\) and \(Y\ge q^3\), a contradiction.
Thus \(n_0=2\) and \(k=8V_0+4\).
Writing \(u=X+2-2Y\), the exact remainder identity is
\begin{equation}
 \eta=4Y^2(Xu+1)+8Y(X+1)+3-4Y.                    \label{eq:compact-p3}
\end{equation}
If \(u\le-1\), this is negative because
\(Y(X-1)>2(X+1)\).  If \(u\ge2\), subtraction of \(k\) leaves at
least \(4Y^2+8Y(X+1)-4Y-1>0\).  Since \(0<\eta<k\), one must have
\(u\in\{0,1\}\), equivalently \(2Y-X\in\{1,2\}\).  But
\(q\mid X,Y\) and \(q\ge32\).  This contradiction proves
\begin{equation}
 p\ge5,\qquad c>c_3:=4\Delta+3.                   \label{eq:compact-p5}
\end{equation}

At \(S=\Delta I\), the adjacent index \(j\) is odd.  Index one gives
\(|V|=1\), already excluded.  Index three gives \(|V|=4Q-3\).
The two possible signs would imply, respectively,
\[
 \mathcal F(b-4\Delta)=c-c_3,\qquad
 \mathcal F(b+4\Delta)=c+c_3.
\]
Both right sides lie strictly between zero and \(\mathcal F\), since
\(0<c-c_3<c+c_3<2c<\mathcal F\).  Thus neither can be a multiple of
\(\mathcal F\).  Consequently \(j\ge5\) and
\[
 |V|\ge16Q^2-20Q+5>8Q^2.
\]
Since \(Q>\mathcal F>2c\), the triangle inequality gives
\[
 |b|>\frac{8Q^2-c}{\mathcal F}
       >8\mathcal F-\tfrac12>4\mathcal F>4g.
\]
Together with \eqref{eq:compact-absolute}, this proves
\begin{equation}
 L=bV=|b||V|>0,\qquad gy<\tfrac38L<L/2.            \label{eq:compact-bvbound}
\end{equation}

For either weight, the new unshifted identity is
\[
 e=b-cT+k-1<b,\qquad \rho<|b|/D,
\]
using \(c>k\) and \(T\ge1\).  If \(B\le-1\), then
\(Wt\le-L-1/h\).  Equation~\eqref{eq:compact-rational} gives
\[
 n=\varepsilon gy-\kappa+Wt+\rho
     <L/2-L+|b|/D<0.
\]
The last inequality uses \(D>2\) and either \(L=V^2\) with
\(|V|>|b|\), or \(L=|b||V|\) with \(|V|>1\).
This contradicts the positive independent input \(n\).  Hence
\begin{equation}
 B\ge0,\qquad C>0.                               \label{eq:compact-transportpos}
\end{equation}

\subsection{Rank recovery and the staged denominator argument}
\label{subsec:compact-staged}

The outer bounds can now be proved without source semantics.  Let
\(G_0=h(q-F)+\alpha\).  As \(q-F=C+\alpha\),
\[
 G_0=q\alpha+(q-1)C,\qquad q\le G_0\le q^2-q,
\]
while \eqref{eq:compact-portbox} gives
\(0<(M_C+qM_F)J<h(1+2q)\).  Therefore
\begin{equation}
 3q<R<q^4-2.                                      \label{eq:compact-outerbound}
\end{equation}
From \eqref{eq:compact-earlybound}, \(c>2R\).  Also
\(c=\psi_A(p)>2p\): \(p\ge3\), and the Pell recurrence gives more
than geometric growth at \(A\ge3\).  Lemma~\ref{lem:compact-rank}
supplies \(m\) with \(c\mid m\), \(m>2p\).
Since \(c^2\mid I\), the strong norm gives \(f^2\equiv1\pmod c\), so
\[
 V\equiv-c\pmod f,\qquad V\equiv-R\pmod c.
\]
Lemma~\ref{lem:compact-index} applies with \(S=\Delta I\), giving
\begin{equation}
 p=R.                                             \label{eq:compact-rankexact}
\end{equation}
In addition, \(T\ge1\) and \(c>2R\) restore \(b\ge c-R>0\) and
\(V>0\).  No later step needs the stronger, generally false inequality
\(b>c\).  The retained mask parity argument yields even \(q\) and odd
\(J\), but the following denominator cancellation does not require that
parity conclusion.\cite{native61,certificate}

Because \(z\) is integral, multiplying
\eqref{eq:compact-rational} by \(hH\) gives
\begin{equation}
 D\mid hHe.                                      \label{eq:compact-partialdiv}
\end{equation}
But \(H=4Y(X+1)-1\equiv-1\pmod q\) and \(h=q-1\), so
\(\gcd(q,hH)=1\).  Since \(q^4\mid D\),
\begin{equation}
 q^4\mid e.                                      \label{eq:compact-smallmodulus}
\end{equation}
The first norm has parameter \(P_0=2XY^2+1\).
The base-one recurrence gives
\(\psi_{P_0}(n_0)\equiv n_0\pmod{P_0-1}\); hence
\begin{equation}
 k\equiv2n_0\pmod{4XY^2}.                        \label{eq:compact-firstcong}
\end{equation}
Equations~\eqref{eq:compact-smallmodulus}--\eqref{eq:compact-firstcong}
imply \(q^4\mid2n_0-R-1\).

Now \(P_0>A\), since
\(P_0-A=Y(2XY-X-1)>0\).  If \(n_0\ge R\), monotonicity in the
positive index and base gives
\[
 k/2=\psi_{P_0}(n_0)\ge\psi_A(R)=c>k,
\]
impossible.  Thus \(1\le n_0<R\), and
\[
 |2n_0-R-1|<R<q^4.
\]
The bounded multiple of \(q^4\) must vanish:
\begin{equation}
 2n_0=R+1.                                       \label{eq:compact-firstindex}
\end{equation}
Only now return to the \emph{full} recurrence modulus in
\eqref{eq:compact-firstcong}.  It gives
\[
 4XY^2\mid e,\qquad D\mid e,\qquad \rho\in\Z.
\]
The middle implication uses the second half of
\eqref{eq:compact-admissible}.  No assumption that \(D\) itself is
coprime to \(hH\) has entered the argument.

\subsection{The retained projection and native calibration}
\label{subsec:compact-transfer}

Write \(R=2r+1\); equation~\eqref{eq:compact-outerbound} gives \(r\ge48\).
The elementary recurrence bounds
\[
 \psi_A(j)\ge(2A-1)^{j-1},\qquad
 \psi_{P_0}(j)\le(2P_0)^{j-1}\qquad(j\ge1)
\]
give a convenient self-contained height estimate.  Indeed,
\[
 (2A-1)^2-2XP_0
 =4Y^2(2X+1)+4Y(X+1)+1-2X>0,
\]
so
\[
 \frac ck=
 \frac{\psi_A(2r+1)}{2\psi_{P_0}(r+1)}
 \ge\frac12\left(\frac{(2A-1)^2}{2P_0}\right)^r
 >\frac{X^r}{2}.
\]
Since \(Y<c/k<Y+1\) and \(X^r>6\), this proves
\begin{equation}
 Y>X^r/3,\qquad X^{r+1}<3Y(X+1)<H.                \label{eq:compact-ratiobound}
\end{equation}

After \(\rho\) is integral, the carrier gives \(H\mid hK_+\).
The needed projection can also be proved directly.  Modulo the odd
integer \(H=4A-5\), put \(u_0=0\) and
\(u_j=3\cdot2^{j-1}\psi_A(j)\) for \(j\ge1\).
The Pell recurrence and \(4A\equiv5\pmod H\) give
\[
 u_{j+2}\equiv5u_{j+1}-4u_j\pmod H,\qquad u_0=0,\quad u_1=3.
\]
The sequence \(4^j-1\) has the same recurrence and initial values.
Thus, writing \(x_*=2^R\),
\[
 3x_*c\equiv2(x_*^2-1)\pmod H,\qquad
 x_*K_+\equiv-2(X-x_*)(x_*X+1)\pmod H.
\]
Multiplying the known divisibility by \(x_*\) and cancelling only the
unit two gives
\begin{equation}
 H\mid h(X-2^R)(2^RX+1).                         \label{eq:compact-projection}
\end{equation}
There is no division by three in this argument.  For the height of the
\emph{entire product}, use \(h<X\), \(X+x_*<Xx_*\), and
\(x_*X+1<2x_*X\):
\[
 \bigl|h(X-x_*)(x_*X+1)\bigr|
 <2x_*^2X^3=8\cdot16^rX^3<X^{r+1}<H.
\]
The penultimate inequality follows from \(X\ge32\) and \(r\ge48\)
(already \(r>13\) is enough).  The product must vanish, and hence
\begin{equation}
 X=2^R.                                          \label{eq:compact-power}
\end{equation}
This is not a cancellation of possibly nonunit individual factors modulo
a composite integer.  Since \(x_*\) is a power of two and \(H\) is odd,
the same congruence now gives \(H\mid K_+\), so \(\kappa\in\Z\).
The carrier then makes \(Wt\) integral; since
\(W\equiv1\pmod h\), \(h\mid B\), and \(t\in\Z_{\ge0}\).

At this point the remaining inherited interface is explicit:
the fixed-minus parity, exact half-binomial rounding, and indexed native
calibration are retained from the supplied ordinary sources.  They are
not inferred from the finite circuit checks.\cite{certificate,native61,indexed}
All hypotheses of that retained fixed-minus parity, exact half-binomial
rounding, and native reconstruction are now available, including the
literal \(q^3\mid Y\).  Reconstruction uses the fresh source base
\(Y(X+1)+2\); it does not identify that base with the present \(A\).
The whole indexed source calibration sends these data to a finite
syntactic ZFC refutation.  This establishes soundness at the explicitly
retained ordinary calibration boundary.

\subsection{Positive completeness}
\label{subsec:compact-completeness}

Conversely, a finite refutation supplies the inherited canonical outer
data, with \(X=2^R\), \(R>3q\), \(R\equiv3\pmod4\),
\(q^3\mid Y\), and an integer \(t=B/h\ge0\).
The constructive scale-one ratio theorem supplies
\begin{align*}
 c&=\psi_A(R),&
 k&=2\psi_{P_0}((R+1)/2),&
 \tau&=\chi_{P_0}((R+1)/2),\\
 \eta&=c-kY>0,&\zeta&=k-\eta>0.
\end{align*}
Here \(P_0=2XY^2+1\).  The first-index congruence gives
\(4XY^2\mid e\), so \eqref{eq:compact-admissible} gives
\(\rho=e/D\in\Z_{\ge0}\).  Projection gives positive integral
\(\kappa\).

The Pell matrix
\[
 \begin{pmatrix}A&\Delta\\1&A\end{pmatrix}
\]
has determinant one, and hence finite order modulo \(c^2\tau\).
A positive multiple of this period supplies a seed pair \((f_0,I_0)\)
with \(c^2\tau\mid I_0\).  Put \(i=I_0/(c^2\tau)>0\), so
\(U=I_0/\tau\).  Then
\[
 a=\Delta I_0^2=f_0^2-1,\qquad
 F_0=f_0^2,\qquad g=\tau f_0.
\]
Doubling this single seed pair gives
\[
 f=2f_0^2-1=1+2a,\qquad
 I=2f_0I_0=2Ug,
\]
exactly the proof pair \eqref{eq:compact-proofpair}.  There is no
simultaneous-solvability assertion about unrelated Pell equations.

The auxiliary witnesses can also be written directly.  With
\(S=\Delta I\) and \(R=2v+1\), where \(v\) is odd, set
\[
 V_*=\frac{\chi_S(R)}{S}=Q_v(S^2),\qquad
 y_* =\psi_S(R).
\]
They are positive integers satisfying the adjacent norm.  Since
\(S^2\equiv1-A^2\pmod{f^2}\) and \(c\mid S\),
Lemma~\ref{lem:compact-odd} gives
\[
 V_*\equiv-c\pmod{f^2},\qquad V_*\equiv-R\pmod c.
\]
The strong norm gives \(f^2\equiv1\pmod c\) and
\(\gcd(c,f^2)=1\).  Therefore
\begin{equation}
 T=\frac{V_*+c+Rf^2}{cf^2}\in\Z_{>0},\qquad
 f^2(cT-R)-c=V_* .                              \label{eq:compact-Tcomplete}
\end{equation}
In particular, no subtraction of \(\alpha\) is needed.

For either prescribed positive weight \(L\), choose
\[
 n=gy_*+(1+hL)t+\rho-\kappa.
\]
All summands are integers.  The inequalities
\[
 f_0>(A-1)I_0>Xc>\kappa
\]
show \(g>\kappa\); the last inequality follows already from
\(H>3\) and \(K_+<3Xc\).  Since \(y_*\ge1\), \(t,\rho\ge0\),
this gives \(n>0\).  Its carrier is exactly \(z=gy_*\), so
\eqref{eq:compact-final} vanishes.  Every independent coordinate is
strictly positive.  This proves completeness by fresh canonical
witnesses and completes the proof of Theorem~\ref{thm:compact-main}.

\subsection{Exact degrees, arithmetic costs, and finite evidence}
\label{subsec:compact-metrics}

Removing the private \(T+\alpha\) addition saves one operation.
Sharing \(XY\) in the main base and first norm saves a multiplication
when \(D=XY\).  Reusing the already computed \(kY\) in
\(\ell=(XY)(kY)\) saves another multiplication.  The \(q^4\) option
then adds the single square \((q^2)^2\).  Counts were obtained from
the complete live graphs, including fixed scalar multiplications:
\begin{center}
\begin{tabular}{@{}llrrrr@{}}
\toprule
\(D\)&\(L\)&Products&Sums&Differences&Total\\
\midrule
\(XY\)&\(V^2\)&44&17&14&75\\
\(XY\)&\(bV\)&45&17&14&76\\
\(q^4\)&\(V^2\)&45&17&14&76\\
\(q^4\)&\(bV\)&46&17&14&77\\
\bottomrule
\end{tabular}
\end{center}

Write \(q_0=\beta J\), \(k_0=\eta+\zeta\).
The full highest homogeneous components, in table order, are
\begin{align}
 &2^{14}q_0^{378}w^{48}s^{110}i^{20}k_0^{64}T^4
       (q_0-F-\alpha)^2,\nonumber\\
 &-2^{10}q_0^{284}w^{36}s^{82}i^{14}k_0^{46}T^2,\nonumber\\
 &2^{14}q_0^{378}w^{46}s^{108}i^{20}k_0^{64}T^4
       (q_0-F-\alpha)^2,\nonumber\\
 &-2^{10}q_0^{284}w^{34}s^{80}i^{14}k_0^{46}T^2.
 \label{eq:compact-leading}
\end{align}
Their degrees are \(626,464,622,460\).  The selected coefficients
\(2^{14}\beta^{378}\) and \(-2^{10}\beta^{284}\) are nonzero at the
actual fixed \(\beta\).  These are complete multivariate leading
components, not numerical rays.  The squared-weight outputs have 390
leading monomials; the \(bV\) outputs have 47.

The supplied standard-library checker proves 32 or 33 exact
free-polynomial cut identities covering every actual gate, checks full
topology and liveness, collects every leading component, and validates
the four literal coefficient payloads and bit lengths.  A separate
reconstruction script regenerates the integers from the finite code in
Section~\ref{sec:compact-constants}, byte for byte.  Normal and optimized
Python runs agree.  A separately tasked AI reviewer checked the ordinary
proof and independently verified the complete leading forms and graph
liveness; its 480 signed modular scalar comparisons are additional
regression evidence, not substitutes for the exact identities or the
unbounded proof.  The independent review led to the explicit folding
proof printed above.

The emitted JSON graphs, coefficient bytes, builders, checkers,
receipts, source ledger, and copied ordinary proof sources are included
in the archive's \texttt{compact/} directory.  They specify actual fixed
polynomials.  The native/refutation theorem remains the expressly
retained calibration; the finite programs do not establish consistency
or replace that theorem.


\section{Verification, comparison, and remaining limits}\label{sec:verification}

The archive separates mathematical proofs from finite artifact checks.
The main conclusions of the unbounded proofs are the centered-residue
theorem, its pre-factorial margin, the compatible odd-square radix
choice, the natural pooling theorem, and the compact admissible-denominator
criterion. Independent review examined the signed branches, order of
recovery, parity choices, and completeness maps.

For the direct quartics, verification reconstructs the exact local source
mutation, checks privacy of every deleted coordinate, replays the monic
projection and integer translation, expands the entire sum of squares,
and recomputes degree, support, height, and variable liveness. Natural
projections additionally verify their nonnegative restoration certificates.
The JSON coefficient arrays are compared exactly; no floating-point
arithmetic is involved.

For the compact equations, exact free-polynomial cut identities cover
every gate of the full binary diagrams. The checker collects their
complete highest homogeneous parts, including all 390 or 47 terms in
the respective leading expressions, and verifies the claimed exact
degrees. A separately implemented audit checks graph topology, liveness,
and the same full leading parts. Its 480 modular comparisons with an
independent scalar implementation are regression checks, not a proof
of the unbounded Pell argument.

All four compact coefficient integers are independently reconstructed
from the published finite 436-character seed and indexed tile formulas.
Their canonical uncompressed bytes agree exactly with the delivered
payloads. This checks the coherent whole tuple, rather than selecting
individually convenient coefficients from different source versions.

The root verification entry point runs the selected checkers in both
ordinary and optimized Python, rebinds the generated report inputs to
the literal equations, and checks the release manifest. Actual receipts
are supplied. Optimized execution ensures that necessary checks have
not inadvertently been expressed only as removable Python assertions.
The PDF was compiled from the shipped LaTeX and inspected for clipping,
broken formulas, table layout, and equation transcription.

None of these finite checks determines whether one of the consistency
equations has a zero. The ordinary proofs, together with the explicitly
identified inherited logical calibration, establish the equivalence.
The delivered package contains no new complete Lean/Rocq derivation of
that endpoint and no claim to have proved the universal nonvanishing
by numerical search.


The new radius theorem is a reusable arithmetic reduction: it applies whenever the surrounding equations identify small canonical residues and a short centered representative is enough. The odd-square step is more specific. It uses the actual odd CRT periods and the freedom to choose an even beta code; it must not be applied to an arbitrary radix without those completeness arguments. Likewise, the compact denominator theorem separates recovery of a small modulus from recovery of the full denominator. Neither reduction is justified by deleting a row and observing that a few evaluations still agree.

Two plausible avenues did not supply a stronger completed result. Coefficient separation can replace prime-factor arguments in a finite Horn flow, but that mechanism already appears in the inspected repository and the current CRT carrier independently uses the factorial. Also, combining unequal natural-domain digit budgets by multiplying one by $q$ yields a correct quadratic constraint, but can obstruct later affine eliminations and increase the final variable count. The retained measured endpoints reflect the complete compiled polynomial, not the local source count.

No global optimum or literature-wide priority is claimed. The repository audit fixes the comparison to the stated commit and its six incoming archives. The most compact new equation still uses the existing native coefficient tuple; the direct equation provides a separate route without that tuple. Further reductions should target the remaining power-certificate structure or the number of independently authenticated digit groups. A claim below the established nine-positive-variable frontier would require substantially different mathematics; this work does not make that claim.

\appendix
\section{Complete residual list of the direct integer quartic}\label{app:direct}

All variables in this appendix range independently over $\Z$. Put
\[
P_{\mathrm{dir}}=\sum_{j=0}^{\SignedRows-1} R_j^2.
\]
The ordered variables are $x_0,\ldots,x_{\SignedVariables-1}$. Every displayed coefficient is literal. The adjacent machine-readable file also supplies the semantic variable names and the fully collected quartic.
% Generated by render_report_inputs.py; do not edit manually.
% Source SHA-256: 36a30a3fb0696a8b02106cc5667b3eb33ef706bd81cc19dbb3d80c69c1846a5e
\begingroup
\footnotesize
\setlength{\jot}{3pt}
\allowdisplaybreaks[4]
\begin{align*}
R_{0} &={}-x_{3}^{2} +x_{4} \\[3pt]
R_{1} &={}x_{0} -x_{0}x_{2} -x_{2} +x_{9} \\[3pt]
R_{2} &={}-1 +x_{1} -x_{2}x_{7} -x_{3} \\[3pt]
R_{3} &={}-1 +x_{1} -x_{8}x_{9} \\[3pt]
R_{4} &={}1 +x_{4} -x_{5}x_{20} +x_{18} \\[3pt]
R_{5} &={}-1 +2x_{10} -x_{10}^{2} -2x_{10}x_{11} +2x_{11} -x_{11}^{2} \\*
&\quad{}+x_{21} \\[3pt]
R_{6} &={}-5 -30x_{10} -36x_{10}^{2} -96x_{10}x_{11} -24x_{10}x_{21} -40x_{11} \\*
&\quad{}-64x_{11}^{2} -32x_{11}x_{21} -10x_{21} -4x_{21}^{2} +x_{22} \\[3pt]
R_{7} &={}-2 +2x_{10} -x_{10}^{2} -4x_{10}x_{22} +4x_{22} -4x_{22}^{2} \\*
&\quad{}+x_{23} \\[3pt]
R_{8} &={}1 +6x_{10} -36x_{10}^{2} -192x_{10}x_{22} -24x_{10}x_{23} +16x_{22} \\*
&\quad{}-256x_{22}^{2} -64x_{22}x_{23} +2x_{23} -4x_{23}^{2} +x_{24} \\[3pt]
R_{9} &={}-2 -x_{11}^{2} -2x_{11}x_{12} -x_{12}^{2} +x_{25} \\[3pt]
R_{10} &={}-9 +7x_{10} -x_{10}^{2} -4x_{10}x_{11} -4x_{10}x_{25} +12x_{11} \\*
&\quad{}-4x_{11}^{2} -8x_{11}x_{25} +12x_{25} -4x_{25}^{2} +x_{26} \\[3pt]
R_{11} &={}-1 +2x_{10} -x_{10}^{2} -2x_{10}x_{12} +2x_{12} -x_{12}^{2} \\*
&\quad{}+x_{27} \\[3pt]
R_{12} &={}-6x_{10} -144x_{10}^{2} -192x_{10}x_{11} -192x_{10}x_{12} -48x_{10}x_{21} -48x_{10}x_{27} \\*
&\quad{}-8x_{11} -64x_{11}^{2} -128x_{11}x_{12} -32x_{11}x_{21} -32x_{11}x_{27} -64x_{12}^{2} \\*
&\quad{}-32x_{12}x_{21} -32x_{12}x_{27} -2x_{21} -4x_{21}^{2} -8x_{21}x_{27} -4x_{27}^{2} \\*
&\quad{}+x_{28} \\[3pt]
R_{13} &={}-4x_{26}^{2} -8x_{26}x_{28} -4x_{28}^{2} +x_{29} \\[3pt]
R_{14} &={}-6 +16x_{12} -16x_{12}^{2} -16x_{12}x_{25} -16x_{12}x_{27} +8x_{25} \\*
&\quad{}-4x_{25}^{2} -8x_{25}x_{27} +8x_{27} -4x_{27}^{2} +x_{30} \\[3pt]
R_{15} &={}1 +6x_{10} -36x_{10}^{2} -96x_{10}x_{11} -48x_{10}x_{12} -24x_{10}x_{21} \\*
&\quad{}-48x_{10}x_{27} -24x_{10}x_{30} +8x_{11} -64x_{11}^{2} -64x_{11}x_{12} -32x_{11}x_{21} \\*
&\quad{}-64x_{11}x_{27} -32x_{11}x_{30} +8x_{12} -16x_{12}^{2} -16x_{12}x_{21} -32x_{12}x_{27} \\*
&\quad{}-16x_{12}x_{30} +2x_{21} -4x_{21}^{2} -16x_{21}x_{27} -8x_{21}x_{30} +8x_{27} \\*
&\quad{}-16x_{27}^{2} -16x_{27}x_{30} +4x_{30} -4x_{30}^{2} +x_{31} \\[3pt]
R_{16} &={}-10 -5x_{10} -x_{10}^{2} -12x_{10}x_{11} -16x_{10}x_{12} -4x_{10}x_{25} \\*
&\quad{}-36x_{11} -36x_{11}^{2} -96x_{11}x_{12} -24x_{11}x_{25} -48x_{12} -64x_{12}^{2} \\*
&\quad{}-32x_{12}x_{25} -12x_{25} -4x_{25}^{2} +x_{32} \\[3pt]
R_{17} &={}-2 -24x_{10} -36x_{10}^{2} -12x_{10}x_{11} -96x_{10}x_{12} -24x_{10}x_{27} \\*
&\quad{}-3x_{11} -x_{11}^{2} -16x_{11}x_{12} -4x_{11}x_{27} -32x_{12} -64x_{12}^{2} \\*
&\quad{}-32x_{12}x_{27} -8x_{27} -4x_{27}^{2} +x_{33} \\[3pt]
R_{18} &={}2 -4x_{32}^{2} -8x_{32}x_{33} -4x_{33}^{2} +x_{34} \\[3pt]
R_{19} &={}-12 +16x_{10} -4x_{10}^{2} -8x_{10}x_{12} -8x_{10}x_{21} -8x_{10}x_{25} \\*
&\quad{}+16x_{12} -4x_{12}^{2} -8x_{12}x_{21} -8x_{12}x_{25} +16x_{21} -4x_{21}^{2} \\*
&\quad{}-8x_{21}x_{25} +16x_{25} -4x_{25}^{2} +x_{35} \\[3pt]
R_{20} &={}-3 -24x_{10} -64x_{10}^{2} -160x_{10}x_{12} -64x_{10}x_{25} -32x_{10}x_{27} \\*
&\quad{}-32x_{10}x_{35} -26x_{12} -100x_{12}^{2} -80x_{12}x_{25} -40x_{12}x_{27} -40x_{12}x_{35} \\*
&\quad{}-8x_{25} -16x_{25}^{2} -16x_{25}x_{27} -16x_{25}x_{35} -6x_{27} -4x_{27}^{2} \\*
&\quad{}-8x_{27}x_{35} -4x_{35} -4x_{35}^{2} +x_{36} \\[3pt]
R_{21} &={}-1 -18x_{10} -36x_{10}^{2} -96x_{10}x_{11} -24x_{10}x_{21} -24x_{11} \\*
&\quad{}-64x_{11}^{2} -32x_{11}x_{21} -6x_{21} -4x_{21}^{2} +x_{37} \\[3pt]
R_{22} &={}-2 +2x_{10} -x_{10}^{2} -4x_{10}x_{37} +4x_{37} -4x_{37}^{2} \\*
&\quad{}+x_{38} \\[3pt]
R_{23} &={}1 +6x_{10} -36x_{10}^{2} -192x_{10}x_{37} -24x_{10}x_{38} +16x_{37} \\*
&\quad{}-256x_{37}^{2} -64x_{37}x_{38} +2x_{38} -4x_{38}^{2} +x_{39} \\[3pt]
R_{24} &={}-x_{11}^{2} -4x_{11}x_{39} -4x_{39}^{2} +x_{40} \\[3pt]
R_{25} &={}-6 -30x_{10} -64x_{10}^{2} -128x_{10}x_{11} -64x_{10}x_{21} -32x_{11} \\*
&\quad{}-64x_{11}^{2} -64x_{11}x_{21} -14x_{21} -16x_{21}^{2} +x_{41} \\[3pt]
R_{26} &={}-4 +8x_{11} -64x_{11}^{2} -96x_{11}x_{12} -32x_{11}x_{25} -32x_{11}x_{41} \\*
&\quad{}+6x_{12} -36x_{12}^{2} -24x_{12}x_{25} -24x_{12}x_{41} +2x_{25} -4x_{25}^{2} \\*
&\quad{}-8x_{25}x_{41} +4x_{41} -4x_{41}^{2} +x_{42} \\[3pt]
R_{27} &={}8x_{10} -64x_{10}^{2} -96x_{10}x_{12} -32x_{10}x_{27} -32x_{10}x_{42} +6x_{12} \\*
&\quad{}-36x_{12}^{2} -24x_{12}x_{27} -24x_{12}x_{42} +2x_{27} -4x_{27}^{2} -8x_{27}x_{42} \\*
&\quad{}+4x_{42} -4x_{42}^{2} +x_{43} \\[3pt]
R_{28} &={}-x_{12}^{2} -2x_{12}x_{13} -x_{13}^{2} +x_{44} \\[3pt]
R_{29} &={}2x_{10} -x_{10}^{2} -4x_{10}x_{12} -4x_{10}x_{44} +4x_{12} -4x_{12}^{2} \\*
&\quad{}-8x_{12}x_{44} +4x_{44} -4x_{44}^{2} +x_{45} \\[3pt]
R_{30} &={}-35 +24x_{10} -4x_{10}^{2} -8x_{10}x_{11} -8x_{10}x_{25} -8x_{10}x_{27} \\*
&\quad{}+24x_{11} -4x_{11}^{2} -8x_{11}x_{25} -8x_{11}x_{27} +24x_{25} -4x_{25}^{2} \\*
&\quad{}-8x_{25}x_{27} +24x_{27} -4x_{27}^{2} +x_{46} \\[3pt]
R_{31} &={}-16x_{10}^{2} -32x_{10}x_{12} -32x_{10}x_{27} -32x_{10}x_{44} -16x_{10}x_{45} -16x_{10}x_{46} \\*
&\quad{}-4x_{12} -16x_{12}^{2} -32x_{12}x_{27} -32x_{12}x_{44} -16x_{12}x_{45} -16x_{12}x_{46} \\*
&\quad{}-16x_{27}^{2} -32x_{27}x_{44} -16x_{27}x_{45} -16x_{27}x_{46} -4x_{44} -16x_{44}^{2} \\*
&\quad{}-16x_{44}x_{45} -16x_{44}x_{46} -2x_{45} -4x_{45}^{2} -8x_{45}x_{46} -4x_{46}^{2} \\*
&\quad{}+x_{47} \\[3pt]
R_{32} &={}x_{12} -25x_{12}^{2} -40x_{12}x_{27} -20x_{12}x_{30} -16x_{27}^{2} -16x_{27}x_{30} \\*
&\quad{}-4x_{30}^{2} +x_{48} \\[3pt]
R_{33} &={}-1 +8x_{11} -4x_{11}^{2} -8x_{11}x_{25} +8x_{25} -4x_{25}^{2} \\*
&\quad{}+x_{49} \\[3pt]
R_{34} &={}-1 -16x_{11}^{2} -8x_{11}x_{12} -32x_{11}x_{25} -16x_{11}x_{49} -x_{12}^{2} \\*
&\quad{}-8x_{12}x_{25} -4x_{12}x_{49} -16x_{25}^{2} -16x_{25}x_{49} -4x_{49}^{2} +x_{50} \\[3pt]
R_{35} &={}1 +x_{11} -x_{11}^{2} -4x_{11}x_{12} -4x_{11}x_{50} -4x_{12}^{2} \\*
&\quad{}-8x_{12}x_{50} -4x_{50}^{2} +x_{51} \\[3pt]
R_{36} &={}-1 -4x_{51} -4x_{51}^{2} +x_{52} \\[3pt]
R_{37} &={}-1 -x_{12}^{2} -8x_{12}x_{51} -4x_{12}x_{52} -16x_{51}^{2} -16x_{51}x_{52} \\*
&\quad{}-4x_{52}^{2} +x_{53} \\[3pt]
R_{38} &={}2 +x_{11} -x_{11}^{2} -4x_{11}x_{12} -4x_{11}x_{53} -4x_{12}^{2} \\*
&\quad{}-8x_{12}x_{53} -4x_{53}^{2} +x_{54} \\[3pt]
R_{39} &={}-4x_{48}^{2} -8x_{48}x_{54} -4x_{54}^{2} +x_{55} \\[3pt]
R_{40} &={}-9 +7x_{10} -x_{10}^{2} -8x_{10}x_{48} -4x_{10}x_{55} +24x_{48} \\*
&\quad{}-16x_{48}^{2} -16x_{48}x_{55} +12x_{55} -4x_{55}^{2} +x_{56} \\[3pt]
R_{41} &={}-1 -4x_{56} -4x_{56}^{2} +x_{57} \\[3pt]
R_{42} &={}-x_{12}^{2} -8x_{12}x_{56} -4x_{12}x_{57} -16x_{56}^{2} -16x_{56}x_{57} -4x_{57}^{2} \\*
&\quad{}+x_{58} \\[3pt]
R_{43} &={}-x_{10} -x_{10}^{2} -4x_{10}x_{12} -4x_{10}x_{58} -4x_{12} -4x_{12}^{2} \\*
&\quad{}-8x_{12}x_{58} -4x_{58} -4x_{58}^{2} +x_{59} \\[3pt]
R_{44} &={}-10 -12x_{59} -4x_{59}^{2} +x_{60} \\[3pt]
R_{45} &={}1 +4x_{11} -4x_{11}^{2} -8x_{11}x_{25} -16x_{11}x_{59} -8x_{11}x_{60} \\*
&\quad{}+4x_{25} -4x_{25}^{2} -16x_{25}x_{59} -8x_{25}x_{60} +8x_{59} -16x_{59}^{2} \\*
&\quad{}-16x_{59}x_{60} +4x_{60} -4x_{60}^{2} +x_{61} \\[3pt]
R_{46} &={}x_{11} -25x_{11}^{2} -40x_{11}x_{25} -20x_{11}x_{61} -16x_{25}^{2} -16x_{25}x_{61} \\*
&\quad{}-4x_{61}^{2} +x_{62} \\[3pt]
R_{47} &={}-10 -12x_{54} -4x_{54}^{2} +x_{63} \\[3pt]
R_{48} &={}1 -16x_{54}^{2} -16x_{54}x_{62} -16x_{54}x_{63} -4x_{62}^{2} -8x_{62}x_{63} \\*
&\quad{}-4x_{63}^{2} +x_{64} \\[3pt]
R_{49} &={}x_{12} -x_{12}^{2} -8x_{12}x_{62} -4x_{12}x_{64} -16x_{62}^{2} -16x_{62}x_{64} \\*
&\quad{}-4x_{64}^{2} +x_{65} \\[3pt]
R_{50} &={}-4x_{10} -4x_{10}^{2} -4x_{10}x_{12} -8x_{10}x_{21} -x_{12} -x_{12}^{2} \\*
&\quad{}-4x_{12}x_{21} -4x_{21} -4x_{21}^{2} +x_{66} \\[3pt]
R_{51} &={}1 +4x_{10} -4x_{10}^{2} -8x_{10}x_{21} -8x_{10}x_{66} +4x_{21} \\*
&\quad{}-4x_{21}^{2} -8x_{21}x_{66} +4x_{66} -4x_{66}^{2} +x_{67} \\[3pt]
R_{52} &={}-2 -4x_{10} -16x_{10}^{2} -32x_{10}x_{21} -16x_{10}x_{65} -16x_{10}x_{67} \\*
&\quad{}-4x_{21} -16x_{21}^{2} -16x_{21}x_{65} -16x_{21}x_{67} -4x_{65}^{2} -8x_{65}x_{67} \\*
&\quad{}-2x_{67} -4x_{67}^{2} +x_{68} \\[3pt]
R_{53} &={}-2 -4x_{68} -4x_{68}^{2} +x_{69} \\[3pt]
R_{54} &={}-x_{3}x_{15} +x_{70} \\[3pt]
R_{55} &={}x_{14} -x_{15} -2x_{24}x_{70} -x_{71} \\[3pt]
R_{56} &={}x_{14} -x_{15} -4x_{26}x_{70} -2x_{29}x_{70} -x_{72} \\[3pt]
R_{57} &={}x_{14} -x_{15} -2x_{31}x_{70} -6x_{70} -x_{73} \\[3pt]
R_{58} &={}x_{14} -x_{15} -4x_{32}x_{70} -2x_{34}x_{70} -x_{74} \\[3pt]
R_{59} &={}x_{14} -x_{15} -2x_{36}x_{70} +6x_{70} -x_{75} \\[3pt]
R_{60} &={}x_{14} -x_{15} -4x_{39}x_{70} -2x_{40}x_{70} -x_{76} \\[3pt]
R_{61} &={}x_{14} -x_{15} -2x_{43}x_{70} -x_{77} \\[3pt]
R_{62} &={}x_{14} -x_{15} -2x_{47}x_{70} -4x_{70} -x_{78} \\[3pt]
R_{63} &={}-x_{10}x_{11} +x_{10}x_{12} +x_{11} +x_{11}x_{12} -x_{12} -x_{12}^{2} \\*
&\quad{}+x_{79} \\[3pt]
R_{64} &={}1 +2x_{13} +x_{13}^{2} -x_{79}^{2} +x_{80} \\[3pt]
R_{65} &={}-x_{80}^{2} +x_{81} \\[3pt]
R_{66} &={}-4x_{68}x_{70} -2x_{69}x_{70} +2x_{70} +x_{82} \\[3pt]
R_{67} &={}-x_{14}^{2} +2x_{14}x_{15} +2x_{14}x_{82} -x_{15}^{2} -2x_{15}x_{82} -x_{82}^{2} \\*
&\quad{}+x_{83} \\[3pt]
R_{68} &={}-1 -x_{3}x_{11} +x_{84} \\[3pt]
R_{69} &={}-x_{15}x_{84} +x_{86}x_{87} +x_{88} \\[3pt]
R_{70} &={}-1 +x_{3} -x_{3}x_{10} +x_{85} \\[3pt]
R_{71} &={}-1 +2x_{3} -2x_{3}x_{10} -2x_{3}x_{21} +x_{86} \\[3pt]
R_{72} &={}-x_{14}x_{85} +x_{87} \\[3pt]
R_{73} &={}-1 -6x_{3}x_{10} -8x_{3}x_{11} -2x_{3}x_{21} +x_{89} \\[3pt]
R_{74} &={}x_{14}x_{84} -x_{15}x_{89} +x_{90} \\[3pt]
R_{75} &={}-x_{14}x_{71} +x_{15}x_{71} +x_{91} \\[3pt]
R_{76} &={}-x_{72}x_{91} +x_{92} \\[3pt]
R_{77} &={}-x_{73}x_{92} +x_{93} \\[3pt]
R_{78} &={}-x_{74}x_{93} +x_{94} \\[3pt]
R_{79} &={}-x_{75}x_{94} +x_{95} \\[3pt]
R_{80} &={}-x_{76}x_{95} +x_{96} \\[3pt]
R_{81} &={}-x_{77}x_{96} +x_{97} \\[3pt]
R_{82} &={}-x_{78}x_{97} +x_{98} \\[3pt]
R_{83} &={}-x_{88}x_{98} +x_{99} \\[3pt]
R_{84} &={}-x_{90}x_{99} +x_{100} \\[3pt]
R_{85} &={}-x_{100}^{2} +x_{101} \\[3pt]
R_{86} &={}-x_{2} -x_{2}x_{18} +x_{18} +x_{102} \\[3pt]
R_{87} &={}-x_{16}x_{102} +x_{103} \\[3pt]
R_{88} &={}-2x_{1}^{2} +2x_{1}x_{14} +2x_{1}x_{15} +2x_{1}x_{103} +2x_{1}x_{104} -x_{14}^{2} \\*
&\quad{}-2x_{14}x_{103} -x_{15}^{2} -2x_{15}x_{104} +x_{19}x_{106} -x_{81}x_{101} -x_{83}x_{101} \\*
&\quad{}-x_{103}^{2} -x_{104}^{2} \\[3pt]
R_{89} &={}-x_{2}x_{17} +x_{17} -x_{17}x_{102} +x_{104} \\[3pt]
R_{90} &={}-9x_{10} -9x_{11} -4x_{21} -2x_{24} -4x_{26} -2x_{29} \\*
&\quad{}-2x_{31} -4x_{32} -2x_{34} -2x_{36} -4x_{39} -2x_{40} \\*
&\quad{}-2x_{43} -2x_{47} -4x_{68} -2x_{69} +x_{105} \\[3pt]
R_{91} &={}-x_{107}x_{110} +x_{110} +x_{112} \\[3pt]
R_{92} &={}-x_{5}x_{109} +x_{108} -x_{109}x_{112} -x_{111} \\[3pt]
R_{93} &={}-x_{19} -x_{19}x_{114} +x_{117} \\[3pt]
R_{94} &={}x_{113} -x_{115}x_{117} -x_{116} \\[3pt]
R_{95} &={}-x_{117}x_{119} +x_{118} -x_{120} \\[3pt]
R_{96} &={}-x_{117}x_{122} +x_{121} -x_{123} \\[3pt]
R_{97} &={}-x_{117}x_{125} +x_{124} -x_{126} \\[3pt]
R_{98} &={}-x_{19}x_{128} +x_{131} \\[3pt]
R_{99} &={}x_{127} -x_{129}x_{131} -x_{130} \\[3pt]
R_{100} &={}-x_{131}x_{133} +x_{132} -x_{134} \\[3pt]
R_{101} &={}-x_{19}x_{136} +x_{139} \\[3pt]
R_{102} &={}x_{135} -x_{137}x_{139} -x_{138} \\[3pt]
R_{103} &={}-x_{139}x_{141} +x_{140} -x_{142} \\[3pt]
R_{104} &={}-x_{19}x_{109} +x_{146} \\[3pt]
R_{105} &={}x_{143} -x_{144}x_{146} -x_{145} \\[3pt]
R_{106} &={}-x_{117}x_{148} +x_{147} -x_{149} \\[3pt]
R_{107} &={}8 +x_{150} -2x_{150}x_{152} -16x_{152} +x_{153} \\[3pt]
R_{108} &={}-x_{150}x_{153} -9x_{153} +x_{154} \\[3pt]
R_{109} &={}2x_{151} +x_{151}^{2} +x_{153}^{2} -x_{154}^{2} \\[3pt]
R_{110} &={}-x_{151}x_{156} -x_{156} +x_{157} \\[3pt]
R_{111} &={}2x_{155} +x_{155}^{2} +x_{156}^{2} -x_{157}^{2} \\[3pt]
R_{112} &={}x_{4}^{2} -2x_{4}x_{151} -2x_{151} +x_{158} \\[3pt]
R_{113} &={}-1 -x_{4}x_{156} +x_{6} -x_{155} -x_{156} +x_{157} \\*
&\quad{}+x_{158}x_{159} \\[3pt]
R_{114} &={}1 -2x_{5} +x_{5}^{2} -2x_{5}x_{151} +x_{160} \\[3pt]
R_{115} &={}-x_{5}x_{206} +x_{19} +x_{160}x_{161} -x_{205} +x_{207} \\[3pt]
R_{116} &={}-1 +2x_{151}x_{221} -x_{162} +2x_{221} -x_{221}^{2} \\[3pt]
R_{117} &={}-1 +x_{107} +x_{162}x_{163} -x_{205} -x_{206}x_{221} +x_{207} \\[3pt]
R_{118} &={}-x_{151}x_{165} -x_{165} +x_{166} \\[3pt]
R_{119} &={}2x_{164} +x_{164}^{2} +x_{165}^{2} -x_{166}^{2} \\[3pt]
R_{120} &={}1 -2x_{107} +x_{107}^{2} -2x_{107}x_{151} +x_{167} \\[3pt]
R_{121} &={}-1 -x_{107}x_{165} +x_{108} -x_{164} +x_{166} +x_{167}x_{168} \\[3pt]
R_{122} &={}4 -4x_{107} +x_{107}^{2} -2x_{107}x_{151} +2x_{151} +x_{169} \\[3pt]
R_{123} &={}-1 -x_{107}x_{156} +x_{109} -x_{155} +x_{156} +x_{157} \\*
&\quad{}+x_{169}x_{170} \\[3pt]
R_{124} &={}1 +x_{151} -x_{151}x_{172} -x_{172} +x_{173} \\[3pt]
R_{125} &={}1 +2x_{171} +x_{171}^{2} -2x_{172} +x_{172}^{2} -x_{173}^{2} \\[3pt]
R_{126} &={}-1 +x_{107} -x_{107}x_{172} +x_{113} +x_{167}x_{174} -x_{171} \\*
&\quad{}+x_{173} \\[3pt]
R_{127} &={}-2x_{151} -x_{151}^{2} +x_{175} \\[3pt]
R_{128} &={}2 -x_{107} -x_{107}x_{155} -x_{107}x_{157} +x_{114} +x_{151}x_{157} \\*
&\quad{}+x_{155} -x_{156}x_{175} +2x_{157} +x_{169}x_{176} \\[3pt]
R_{129} &={}-x_{151}x_{178} -x_{178} +x_{179} \\[3pt]
R_{130} &={}2x_{177} +x_{177}^{2} +x_{178}^{2} -x_{179}^{2} \\[3pt]
R_{131} &={}-1 -x_{107}x_{178} +x_{118} +x_{167}x_{180} -x_{177} +x_{179} \\[3pt]
R_{132} &={}-x_{151}x_{182} -x_{182} +x_{183} \\[3pt]
R_{133} &={}2x_{181} +x_{181}^{2} +x_{182}^{2} -x_{183}^{2} \\[3pt]
R_{134} &={}-1 -x_{107}x_{182} +x_{121} +x_{167}x_{184} -x_{181} +x_{183} \\[3pt]
R_{135} &={}-x_{151}x_{186} -x_{186} +x_{187} \\[3pt]
R_{136} &={}2x_{185} +x_{185}^{2} +x_{186}^{2} -x_{187}^{2} \\[3pt]
R_{137} &={}-1 -x_{107}x_{186} +x_{124} +x_{167}x_{188} -x_{185} +x_{187} \\[3pt]
R_{138} &={}-x_{151}x_{190} -x_{190} +x_{191} \\[3pt]
R_{139} &={}2x_{189} +x_{189}^{2} +x_{190}^{2} -x_{191}^{2} \\[3pt]
R_{140} &={}-1 -x_{107}x_{190} +x_{127} +x_{167}x_{192} -x_{189} +x_{191} \\[3pt]
R_{141} &={}-x_{151}x_{194} -x_{194} +x_{195} \\[3pt]
R_{142} &={}2x_{193} +x_{193}^{2} +x_{194}^{2} -x_{195}^{2} \\[3pt]
R_{143} &={}-1 -x_{107}x_{194} +x_{128} +x_{169}x_{196} -x_{193} +x_{194} \\*
&\quad{}+x_{195} \\[3pt]
R_{144} &={}-x_{151}x_{198} -x_{198} +x_{199} \\[3pt]
R_{145} &={}2x_{197} +x_{197}^{2} +x_{198}^{2} -x_{199}^{2} \\[3pt]
R_{146} &={}-1 -x_{107}x_{198} +x_{132} +x_{167}x_{200} -x_{197} +x_{199} \\[3pt]
R_{147} &={}-x_{151}x_{202} -x_{202} +x_{203} \\[3pt]
R_{148} &={}2x_{201} +x_{201}^{2} +x_{202}^{2} -x_{203}^{2} \\[3pt]
R_{149} &={}-1 -x_{107}x_{202} +x_{135} +x_{167}x_{204} -x_{201} +x_{203} \\[3pt]
R_{150} &={}1 -x_{151}x_{206} -x_{206} +x_{207} \\[3pt]
R_{151} &={}-1 +2x_{205} +x_{205}^{2} +x_{206}^{2} -2x_{207} -x_{207}^{2} \\[3pt]
R_{152} &={}-x_{107}x_{206} +x_{136} +x_{169}x_{208} -x_{205} +x_{206} +x_{207} \\[3pt]
R_{153} &={}-x_{151}x_{210} -x_{210} +x_{211} \\[3pt]
R_{154} &={}2x_{209} +x_{209}^{2} +x_{210}^{2} -x_{211}^{2} \\[3pt]
R_{155} &={}-1 -x_{107}x_{210} +x_{140} +x_{167}x_{212} -x_{209} +x_{211} \\[3pt]
R_{156} &={}-x_{151}x_{214} -x_{214} +x_{215} \\[3pt]
R_{157} &={}2x_{213} +x_{213}^{2} +x_{214}^{2} -x_{215}^{2} \\[3pt]
R_{158} &={}-1 -x_{107}x_{214} +x_{143} +x_{167}x_{216} -x_{213} +x_{215} \\[3pt]
R_{159} &={}-x_{151}x_{218} -x_{218} +x_{219} \\[3pt]
R_{160} &={}2x_{217} +x_{217}^{2} +x_{218}^{2} -x_{219}^{2} \\[3pt]
R_{161} &={}-1 -x_{107}x_{218} +x_{147} +x_{167}x_{220} -x_{217} +x_{219} \\[3pt]
R_{162} &={}1 -2x_{0} -3x_{0}^{2} -2x_{0}x_{4} -x_{1}^{2} -x_{2}^{2} \\*
&\quad{}-2x_{4} -2x_{4}^{2} -x_{5}^{2} -x_{6}^{2} +2x_{10} -x_{10}^{2} \\*
&\quad{}-x_{11}^{2} -x_{12}^{2} -x_{13}^{2} -x_{14}^{2} -x_{15}^{2} -x_{16}^{2} \\*
&\quad{}-x_{17}^{2} -x_{18}^{2} -x_{19}^{2} -x_{105}^{2} +2x_{107} -2x_{107}^{2} \\*
&\quad{}-x_{108}^{2} -x_{109}^{2} -x_{113}^{2} -2x_{114} -x_{114}^{2} -x_{118}^{2} \\*
&\quad{}-x_{121}^{2} -x_{124}^{2} -x_{127}^{2} -x_{128}^{2} -x_{132}^{2} -x_{135}^{2} \\*
&\quad{}-x_{136}^{2} -x_{140}^{2} -x_{143}^{2} -x_{147}^{2} +x_{150} -x_{221}^{2} \\[3pt]
R_{163} &={}-x_{5}x_{10} -x_{5}x_{11} -x_{5}x_{12} -x_{5}x_{13} -x_{5}x_{14} -x_{5}x_{15} \\*
&\quad{}-x_{5}x_{16} -x_{5}x_{17} -x_{5}x_{18} +x_{221} \\[3pt]
R_{164} &={}-16 +x_{0} -x_{222}^{2} \\[3pt]
R_{165} &={}-1 +x_{2} +x_{6} -x_{6}^{2} \\[3pt]
R_{166} &={}-x_{151}x_{224} -x_{224} +x_{233} \\[3pt]
R_{167} &={}-1 +x_{223}^{2} +x_{224}^{2} -x_{233}^{2} \\[3pt]
R_{168} &={}-x_{226} -x_{226}x_{227} +x_{234} \\[3pt]
R_{169} &={}-1 +x_{225}^{2} +x_{226}^{2} -x_{234}^{2} \\[3pt]
R_{170} &={}-x_{218}^{2} +x_{235} \\[3pt]
R_{171} &={}x_{224} -2x_{229}x_{235} -x_{235} \\[3pt]
R_{172} &={}-4x_{224} -8x_{224}x_{228} +x_{227} \\[3pt]
R_{173} &={}-x_{151} -x_{223}x_{230} +x_{227} \\[3pt]
R_{174} &={}-1 -x_{217} -x_{223}x_{231} +x_{225} \\[3pt]
R_{175} &={}-x_{105} -4x_{218}x_{232} +x_{226} \\[3pt]
R_{176} &={}2x_{0}^{2} +2x_{0}x_{4} -2x_{0}x_{156} -2x_{0}x_{165} +x_{1}^{2} -2x_{1}x_{206} \\*
&\quad{}+x_{2}^{2} -2x_{2}x_{194} +x_{4}^{2} -2x_{4}x_{165} +x_{10}^{2} -2x_{10}x_{172} \\*
&\quad{}+x_{11}^{2} -2x_{11}x_{178} +x_{12}^{2} -2x_{12}x_{182} +x_{13}^{2} -2x_{13}x_{186} \\*
&\quad{}+x_{14}^{2} -2x_{14}x_{190} +x_{15}^{2} -2x_{15}x_{198} +x_{16}^{2} -2x_{16}x_{202} \\*
&\quad{}+x_{17}^{2} -2x_{17}x_{210} +x_{18}^{2} -2x_{18}x_{214} +x_{105}^{2} -2x_{105}x_{218} \\*
&\quad{}-x_{151}x_{236} +x_{155}^{2} +x_{156}^{2} +x_{164}^{2} +x_{165}^{2} +x_{171}^{2} \\*
&\quad{}+x_{172}^{2} +x_{177}^{2} +x_{178}^{2} +x_{181}^{2} +x_{182}^{2} +x_{185}^{2} \\*
&\quad{}+x_{186}^{2} +x_{189}^{2} +x_{190}^{2} +x_{193}^{2} +x_{194}^{2} +x_{197}^{2} \\*
&\quad{}+x_{198}^{2} +x_{201}^{2} +x_{202}^{2} +x_{205}^{2} +x_{206}^{2} +x_{209}^{2} \\*
&\quad{}+x_{210}^{2} +x_{213}^{2} +x_{214}^{2} +x_{217}^{2} +x_{218}^{2} \\[3pt]
R_{177} &={}-1 -x_{3}x_{5} +x_{5} +x_{6} -x_{237}^{2} -x_{238}^{2} \\*
&\quad{}-x_{239} -x_{239}^{2} \\[3pt]
R_{178} &={}x_{3}x_{5} -x_{6} -x_{240}^{2} -x_{241}^{2} -x_{242} -x_{242}^{2} \\[3pt]
R_{179} &={}-2 -x_{105}^{2} -x_{156}^{2} -x_{165}^{2} +2x_{172} -x_{172}^{2} \\*
&\quad{}-x_{178}^{2} -x_{182}^{2} -x_{186}^{2} -x_{190}^{2} -x_{194}^{2} -x_{198}^{2} \\*
&\quad{}-x_{202}^{2} -x_{206}^{2} -x_{210}^{2} -x_{214}^{2} +x_{218}^{2} -x_{243}^{2} \\*
&\quad{}-x_{244}^{2} -x_{245} -x_{245}^{2} \\[3pt]
R_{180} &={}-2 +x_{107} -4x_{246} -4x_{246}^{2} \\[3pt]
R_{181} &={}-x_{107}x_{111} +x_{111} +x_{247} \\[3pt]
R_{182} &={}-x_{107}x_{145} +x_{145} +x_{248} \\[3pt]
R_{183} &={}2x_{114} +x_{114}^{2} -4x_{116}^{2} -4x_{120}^{2} -4x_{123}^{2} -4x_{126}^{2} \\*
&\quad{}-4x_{149}^{2} -4x_{247}^{2} -4x_{248}^{2} -x_{249}^{2} -x_{250}^{2} -x_{251} \\*
&\quad{}-x_{251}^{2} \\[3pt]
R_{184} &={}-1 +x_{128}^{2} -4x_{130}^{2} -4x_{134}^{2} -x_{252}^{2} -x_{253}^{2} \\*
&\quad{}-x_{254} -x_{254}^{2} \\[3pt]
R_{185} &={}-1 +x_{136}^{2} -4x_{138}^{2} -4x_{142}^{2} -x_{255}^{2} -x_{256}^{2} \\*
&\quad{}-x_{257} -x_{257}^{2}
\end{align*}
\endgroup


\section{Complete residual list of the balanced natural companion}\label{app:natural}

All variables here range over $\N$, including zero. The companion is
$P_{\mathrm{nat}}(x_0,\ldots,x_{221})=\sum_{j=0}^{171}S_j^2$,
with the following literal residuals in its own independent coordinate list.
% Generated by render_report_inputs.py; do not edit manually.
% Source SHA-256: 66fd9c31abd7f071308c1b39b05ebfa97a8185e2eefe77028300a59f8a2f54c5
\begingroup
\footnotesize
\setlength{\jot}{3pt}
\allowdisplaybreaks[4]
\begin{align*}
S_{0} &={}-x_{2}^{2} +x_{3} \\[3pt]
S_{1} &={}-x_{2} -x_{2}x_{5} -x_{2}x_{6} +x_{4} +x_{6} \\[3pt]
S_{2} &={}-13 -x_{0} -x_{0}x_{2} -x_{0}x_{8} -13x_{2} -13x_{8} \\*
&\quad{}+x_{10} \\[3pt]
S_{3} &={}-1 +x_{1} -x_{2} -x_{2}x_{7} -2x_{7} -x_{7}x_{8} \\[3pt]
S_{4} &={}-1 +x_{1} -x_{9} -x_{9}x_{10} \\[3pt]
S_{5} &={}2 +x_{3} -x_{5}x_{21} -x_{6}x_{21} +x_{19} -x_{21} \\[3pt]
S_{6} &={}-6 -5x_{11} -x_{11}^{2} -2x_{11}x_{12} -5x_{12} -x_{12}^{2} \\*
&\quad{}+x_{22} \\[3pt]
S_{7} &={}-72 -68x_{12} -16x_{12}^{2} -16x_{12}x_{22} -34x_{22} -4x_{22}^{2} \\*
&\quad{}+x_{23} \\[3pt]
S_{8} &={}-30 -11x_{11} -x_{11}^{2} -4x_{11}x_{23} -22x_{23} -4x_{23}^{2} \\*
&\quad{}+x_{24} \\[3pt]
S_{9} &={}-380 -312x_{23} -64x_{23}^{2} -32x_{23}x_{24} -78x_{24} -4x_{24}^{2} \\*
&\quad{}+x_{25} \\[3pt]
S_{10} &={}-6 -5x_{12} -x_{12}^{2} -2x_{12}x_{13} -5x_{13} -x_{13}^{2} \\*
&\quad{}+x_{26} \\[3pt]
S_{11} &={}-6 -5x_{11} -x_{11}^{2} -2x_{11}x_{13} -5x_{13} -x_{13}^{2} \\*
&\quad{}+x_{27} \\[3pt]
S_{12} &={}-414 -22x_{11} -2x_{11}^{2} -16x_{11}x_{13} -8x_{11}x_{26} -200x_{12} \\*
&\quad{}-32x_{12}^{2} -64x_{12}x_{13} -32x_{12}x_{22} -32x_{12}x_{27} -320x_{13} -64x_{13}^{2} \\*
&\quad{}-32x_{13}x_{22} -32x_{13}x_{26} -32x_{13}x_{27} -100x_{22} -8x_{22}^{2} -16x_{22}x_{27} \\*
&\quad{}-52x_{26} -8x_{26}^{2} -108x_{27} -8x_{27}^{2} +x_{28} \\[3pt]
S_{13} &={}-1344 -800x_{12} -128x_{12}^{2} -256x_{12}x_{13} -128x_{12}x_{22} -128x_{12}x_{27} \\*
&\quad{}-864x_{13} -128x_{13}^{2} -128x_{13}x_{22} -128x_{13}x_{27} -400x_{22} -32x_{22}^{2} \\*
&\quad{}-64x_{22}x_{27} -432x_{27} -32x_{27}^{2} -2x_{28} -2x_{28}^{2} +x_{29} \\[3pt]
S_{14} &={}-247 -168x_{11} -32x_{11}^{2} -64x_{11}x_{12} -32x_{11}x_{26} -32x_{11}x_{27} \\*
&\quad{}-188x_{12} -32x_{12}^{2} -32x_{12}x_{26} -32x_{12}x_{27} -2x_{22} -92x_{26} \\*
&\quad{}-8x_{26}^{2} -16x_{26}x_{27} -84x_{27} -8x_{27}^{2} +x_{30} \\[3pt]
S_{15} &={}-960 -672x_{11} -128x_{11}^{2} -256x_{11}x_{12} -128x_{11}x_{26} -128x_{11}x_{27} \\*
&\quad{}-736x_{12} -128x_{12}^{2} -128x_{12}x_{26} -128x_{12}x_{27} -368x_{26} -32x_{26}^{2} \\*
&\quad{}-64x_{26}x_{27} -336x_{27} -32x_{27}^{2} -2x_{30} -2x_{30}^{2} +x_{31} \\[3pt]
S_{16} &={}-276 -30x_{11} -2x_{11}^{2} -16x_{11}x_{13} -8x_{11}x_{26} -30x_{12} \\*
&\quad{}-2x_{12}^{2} -16x_{12}x_{13} -8x_{12}x_{27} -272x_{13} -64x_{13}^{2} -32x_{13}x_{26} \\*
&\quad{}-32x_{13}x_{27} -68x_{26} -8x_{26}^{2} -68x_{27} -8x_{27}^{2} +x_{32} \\[3pt]
S_{17} &={}-552 -120x_{12} -8x_{12}^{2} -64x_{12}x_{13} -32x_{12}x_{27} -544x_{13} \\*
&\quad{}-128x_{13}^{2} -128x_{13}x_{27} -272x_{27} -32x_{27}^{2} -2x_{32} -2x_{32}^{2} \\*
&\quad{}+x_{33} \\[3pt]
S_{18} &={}-247 -4x_{11} -352x_{12} -128x_{12}^{2} -64x_{12}x_{22} -64x_{12}x_{26} \\*
&\quad{}-92x_{22} -8x_{22}^{2} -16x_{22}x_{26} -84x_{26} -8x_{26}^{2} -2x_{27} \\*
&\quad{}+x_{34} \\[3pt]
S_{19} &={}-960 -1408x_{12} -512x_{12}^{2} -256x_{12}x_{22} -256x_{12}x_{26} -368x_{22} \\*
&\quad{}-32x_{22}^{2} -64x_{22}x_{26} -336x_{26} -32x_{26}^{2} -2x_{34} -2x_{34}^{2} \\*
&\quad{}+x_{35} \\[3pt]
S_{20} &={}-56 -60x_{12} -16x_{12}^{2} -16x_{12}x_{22} -30x_{22} -4x_{22}^{2} \\*
&\quad{}+x_{36} \\[3pt]
S_{21} &={}-30 -11x_{11} -x_{11}^{2} -4x_{11}x_{36} -22x_{36} -4x_{36}^{2} \\*
&\quad{}+x_{37} \\[3pt]
S_{22} &={}-380 -312x_{36} -64x_{36}^{2} -32x_{36}x_{37} -78x_{37} -4x_{37}^{2} \\*
&\quad{}+x_{38} \\[3pt]
S_{23} &={}-30 -11x_{12} -x_{12}^{2} -4x_{12}x_{38} -22x_{38} -4x_{38}^{2} \\*
&\quad{}+x_{39} \\[3pt]
S_{24} &={}-349 -420x_{12} -128x_{12}^{2} -128x_{12}x_{22} -208x_{22} -32x_{22}^{2} \\*
&\quad{}-2x_{26} +x_{40} \\[3pt]
S_{25} &={}-1368 -1664x_{12} -512x_{12}^{2} -512x_{12}x_{22} -832x_{22} -128x_{22}^{2} \\*
&\quad{}-2x_{40} -2x_{40}^{2} +x_{41} \\[3pt]
S_{26} &={}-56 -60x_{11} -16x_{11}^{2} -16x_{11}x_{27} -8x_{11}x_{41} -30x_{27} \\*
&\quad{}-4x_{27}^{2} -4x_{27}x_{41} -15x_{41} -x_{41}^{2} +x_{42} \\[3pt]
S_{27} &={}-6 -5x_{13} -x_{13}^{2} -2x_{13}x_{14} -5x_{14} -x_{14}^{2} \\*
&\quad{}+x_{43} \\[3pt]
S_{28} &={}-224 -26x_{11} -2x_{11}^{2} -16x_{11}x_{14} -8x_{11}x_{43} -288x_{13} \\*
&\quad{}-128x_{13}^{2} -64x_{13}x_{26} -64x_{13}x_{27} -88x_{14} -32x_{14}^{2} -32x_{14}x_{43} \\*
&\quad{}-76x_{26} -8x_{26}^{2} -16x_{26}x_{27} -68x_{27} -8x_{27}^{2} -44x_{43} \\*
&\quad{}-8x_{43}^{2} +x_{44} \\[3pt]
S_{29} &={}-640 -1152x_{13} -512x_{13}^{2} -256x_{13}x_{26} -256x_{13}x_{27} -304x_{26} \\*
&\quad{}-32x_{26}^{2} -64x_{26}x_{27} -272x_{27} -32x_{27}^{2} -2x_{44} -2x_{44}^{2} \\*
&\quad{}+x_{45} \\[3pt]
S_{30} &={}-241 -168x_{11} -32x_{11}^{2} -64x_{11}x_{12} -32x_{11}x_{26} -32x_{11}x_{27} \\*
&\quad{}-184x_{12} -32x_{12}^{2} -32x_{12}x_{26} -32x_{12}x_{27} -x_{13} -92x_{26} \\*
&\quad{}-8x_{26}^{2} -16x_{26}x_{27} -84x_{27} -8x_{27}^{2} +x_{46} \\[3pt]
S_{31} &={}-42 -52x_{13} -16x_{13}^{2} -16x_{13}x_{26} -26x_{26} -4x_{26}^{2} \\*
&\quad{}+x_{47} \\[3pt]
S_{32} &={}-30 -11x_{13} -x_{13}^{2} -4x_{13}x_{47} -22x_{47} -4x_{47}^{2} \\*
&\quad{}+x_{48} \\[3pt]
S_{33} &={}-306 -35x_{12} -x_{12}^{2} -16x_{12}x_{47} -4x_{12}x_{48} -280x_{47} \\*
&\quad{}-64x_{47}^{2} -32x_{47}x_{48} -70x_{48} -4x_{48}^{2} +x_{49} \\[3pt]
S_{34} &={}-4556 -4320x_{47} -1024x_{47}^{2} -512x_{47}x_{48} -128x_{47}x_{49} -1080x_{48} \\*
&\quad{}-64x_{48}^{2} -32x_{48}x_{49} -270x_{49} -4x_{49}^{2} +x_{50} \\[3pt]
S_{35} &={}-30 -11x_{13} -x_{13}^{2} -4x_{13}x_{50} -22x_{50} -4x_{50}^{2} \\*
&\quad{}+x_{51} \\[3pt]
S_{36} &={}-306 -35x_{12} -x_{12}^{2} -16x_{12}x_{50} -4x_{12}x_{51} -280x_{50} \\*
&\quad{}-64x_{50}^{2} -32x_{50}x_{51} -70x_{51} -4x_{51}^{2} +x_{52} \\[3pt]
S_{37} &={}-1028 -672x_{11} -128x_{11}^{2} -256x_{11}x_{12} -128x_{11}x_{26} -128x_{11}x_{27} \\*
&\quad{}-736x_{12} -128x_{12}^{2} -128x_{12}x_{26} -128x_{12}x_{27} -368x_{26} -32x_{26}^{2} \\*
&\quad{}-64x_{26}x_{27} -336x_{27} -32x_{27}^{2} -2x_{46} -2x_{46}^{2} -32x_{50} \\*
&\quad{}-8x_{51} -2x_{52} +x_{53} \\[3pt]
S_{38} &={}-264 -128x_{50} -32x_{51} -8x_{52} -2x_{53} -2x_{53}^{2} \\*
&\quad{}+x_{54} \\[3pt]
S_{39} &={}-2 -3x_{11} -x_{11}^{2} -2x_{11}x_{54} -3x_{54} -x_{54}^{2} \\*
&\quad{}+x_{55} \\[3pt]
S_{40} &={}-12 -28x_{54} -16x_{54}^{2} -16x_{54}x_{55} -14x_{55} -4x_{55}^{2} \\*
&\quad{}+x_{56} \\[3pt]
S_{41} &={}-30 -11x_{13} -x_{13}^{2} -4x_{13}x_{56} -22x_{56} -4x_{56}^{2} \\*
&\quad{}+x_{57} \\[3pt]
S_{42} &={}-306 -35x_{11} -x_{11}^{2} -16x_{11}x_{56} -4x_{11}x_{57} -280x_{56} \\*
&\quad{}-64x_{56}^{2} -32x_{56}x_{57} -70x_{57} -4x_{57}^{2} +x_{58} \\[3pt]
S_{43} &={}-4556 -4320x_{56} -1024x_{56}^{2} -512x_{56}x_{57} -128x_{56}x_{58} -1080x_{57} \\*
&\quad{}-64x_{57}^{2} -32x_{57}x_{58} -270x_{58} -4x_{58}^{2} +x_{59} \\[3pt]
S_{44} &={}-90 -76x_{13} -16x_{13}^{2} -16x_{13}x_{26} -16x_{13}x_{59} -38x_{26} \\*
&\quad{}-4x_{26}^{2} -8x_{26}x_{59} -38x_{59} -4x_{59}^{2} +x_{60} \\[3pt]
S_{45} &={}-4556 -4320x_{50} -1024x_{50}^{2} -512x_{50}x_{51} -128x_{50}x_{52} -1080x_{51} \\*
&\quad{}-64x_{51}^{2} -32x_{51}x_{52} -270x_{52} -4x_{52}^{2} +x_{61} \\[3pt]
S_{46} &={}-682 -70x_{12} -2x_{12}^{2} -32x_{12}x_{59} -8x_{12}x_{60} -592x_{59} \\*
&\quad{}-128x_{59}^{2} -64x_{59}x_{60} -148x_{60} -8x_{60}^{2} -2x_{61} +x_{62} \\[3pt]
S_{47} &={}-x_{62} -x_{62}^{2} +x_{63} \\[3pt]
S_{48} &={}-42 -52x_{12} -16x_{12}^{2} -8x_{12}x_{13} -16x_{12}x_{22} -13x_{13} \\*
&\quad{}-x_{13}^{2} -4x_{13}x_{22} -26x_{22} -4x_{22}^{2} +x_{64} \\[3pt]
S_{49} &={}-2278 -2264x_{12} -800x_{12}^{2} -800x_{12}x_{22} -160x_{12}x_{64} -70x_{13} \\*
&\quad{}-2x_{13}^{2} -32x_{13}x_{61} -8x_{13}x_{63} -1132x_{22} -200x_{22}^{2} -80x_{22}x_{64} \\*
&\quad{}-592x_{61} -128x_{61}^{2} -64x_{61}x_{63} -148x_{63} -8x_{63}^{2} -228x_{64} \\*
&\quad{}-8x_{64}^{2} +x_{65} \\[3pt]
S_{50} &={}-2712 -280x_{13} -8x_{13}^{2} -128x_{13}x_{61} -32x_{13}x_{63} -2368x_{61} \\*
&\quad{}-512x_{61}^{2} -256x_{61}x_{63} -592x_{63} -32x_{63}^{2} -2x_{65} -2x_{65}^{2} \\*
&\quad{}+x_{66} \\[3pt]
S_{51} &={}-2 -3x_{66} -x_{66}^{2} +x_{67} \\[3pt]
S_{52} &={}-x_{2} -x_{2}x_{16} +x_{68} \\[3pt]
S_{53} &={}-x_{15}^{2} +2x_{15}x_{16} -x_{16}^{2} +x_{69} \\[3pt]
S_{54} &={}-2x_{25}x_{68} -4x_{68} +x_{70} \\[3pt]
S_{55} &={}-x_{15}^{2} +2x_{15}x_{16} +2x_{15}x_{70} -x_{16}^{2} -2x_{16}x_{70} -x_{70}^{2} \\*
&\quad{}+x_{71} \\[3pt]
S_{56} &={}-x_{29}x_{68} +x_{72} \\[3pt]
S_{57} &={}-x_{15}^{2} +2x_{15}x_{16} +2x_{15}x_{72} -x_{16}^{2} -2x_{16}x_{72} -x_{72}^{2} \\*
&\quad{}+x_{73} \\[3pt]
S_{58} &={}-x_{31}x_{68} +x_{74} \\[3pt]
S_{59} &={}-x_{15}^{2} +2x_{15}x_{16} +2x_{15}x_{74} -x_{16}^{2} -2x_{16}x_{74} -x_{74}^{2} \\*
&\quad{}+x_{75} \\[3pt]
S_{60} &={}-x_{33}x_{68} +x_{76} \\[3pt]
S_{61} &={}-x_{15}^{2} +2x_{15}x_{16} +2x_{15}x_{76} -x_{16}^{2} -2x_{16}x_{76} -x_{76}^{2} \\*
&\quad{}+x_{77} \\[3pt]
S_{62} &={}-x_{35}x_{68} +x_{78} \\[3pt]
S_{63} &={}-x_{15}^{2} +2x_{15}x_{16} +2x_{15}x_{78} -x_{16}^{2} -2x_{16}x_{78} -x_{78}^{2} \\*
&\quad{}+x_{79} \\[3pt]
S_{64} &={}-4x_{12}x_{68} -2x_{39}x_{68} -4x_{68} +x_{80} \\[3pt]
S_{65} &={}-x_{15}^{2} +2x_{15}x_{16} +2x_{15}x_{80} -x_{16}^{2} -2x_{16}x_{80} -x_{80}^{2} \\*
&\quad{}+x_{81} \\[3pt]
S_{66} &={}-4x_{41}x_{68} -2x_{42}x_{68} +x_{82} \\[3pt]
S_{67} &={}-x_{15}^{2} +2x_{15}x_{16} +2x_{15}x_{82} -x_{16}^{2} -2x_{16}x_{82} -x_{82}^{2} \\*
&\quad{}+x_{83} \\[3pt]
S_{68} &={}-x_{45}x_{68} +x_{84} \\[3pt]
S_{69} &={}-x_{15}^{2} +2x_{15}x_{16} +2x_{15}x_{84} -x_{16}^{2} -2x_{16}x_{84} -x_{84}^{2} \\*
&\quad{}+x_{85} \\[3pt]
S_{70} &={}-x_{11}^{2} +2x_{11}x_{13} -x_{13}^{2} +x_{86} \\[3pt]
S_{71} &={}-x_{12}^{2} +2x_{12}x_{13} -x_{13}^{2} +x_{87} \\[3pt]
S_{72} &={}-x_{86}x_{87} +x_{88} \\[3pt]
S_{73} &={}-4 -4x_{14} -x_{14}^{2} +x_{89} \\[3pt]
S_{74} &={}-x_{88}^{2} +2x_{88}x_{89} -x_{89}^{2} +x_{90} \\[3pt]
S_{75} &={}-2x_{67}x_{68} -4x_{68} +x_{91} \\[3pt]
S_{76} &={}-x_{15}^{2} +2x_{15}x_{16} +2x_{15}x_{91} -x_{16}^{2} -2x_{16}x_{91} -x_{91}^{2} \\*
&\quad{}+x_{92} \\[3pt]
S_{77} &={}-x_{2} -x_{2}x_{12} +x_{93} \\[3pt]
S_{78} &={}-1 -x_{16} -x_{16}x_{93} -x_{93} +x_{94} \\[3pt]
S_{79} &={}-x_{2} -x_{2}x_{11} +x_{95} \\[3pt]
S_{80} &={}-4x_{2} -4x_{2}x_{12} -2x_{2}x_{22} +x_{96} \\[3pt]
S_{81} &={}-1 -x_{15} -x_{15}x_{95} -x_{95} +x_{97} \\[3pt]
S_{82} &={}-x_{96}x_{97} -x_{97} +x_{98} \\[3pt]
S_{83} &={}-x_{94}^{2} +2x_{94}x_{98} -x_{98}^{2} +x_{99} \\[3pt]
S_{84} &={}-6x_{2} -4x_{2}x_{12} -2x_{2}x_{22} +x_{100} \\[3pt]
S_{85} &={}-1 -x_{16} -x_{16}x_{100} -x_{100} +x_{101} \\[3pt]
S_{86} &={}-1 -x_{15} -x_{15}x_{93} -x_{93} +x_{102} \\[3pt]
S_{87} &={}-x_{101}^{2} +2x_{101}x_{102} -x_{102}^{2} +x_{103} \\[3pt]
S_{88} &={}-x_{69}x_{71} +x_{104} \\[3pt]
S_{89} &={}-x_{73}x_{104} +x_{105} \\[3pt]
S_{90} &={}-x_{75}x_{105} +x_{106} \\[3pt]
S_{91} &={}-x_{77}x_{106} +x_{107} \\[3pt]
S_{92} &={}-x_{79}x_{107} +x_{108} \\[3pt]
S_{93} &={}-x_{81}x_{108} +x_{109} \\[3pt]
S_{94} &={}-x_{83}x_{109} +x_{110} \\[3pt]
S_{95} &={}-x_{85}x_{110} +x_{111} \\[3pt]
S_{96} &={}-x_{90}x_{111} -x_{92}x_{111} +x_{112} \\[3pt]
S_{97} &={}-x_{99}x_{112} +x_{113} \\[3pt]
S_{98} &={}-2 -2x_{2} -x_{2}x_{19} -2x_{8} -x_{8}x_{19} -x_{19} \\*
&\quad{}+x_{114} \\[3pt]
S_{99} &={}-1 -x_{17} -x_{17}x_{114} -x_{114} +x_{115} \\[3pt]
S_{100} &={}-2 -x_{2} -x_{2}x_{18} -x_{8} -x_{8}x_{18} -2x_{18} \\*
&\quad{}-x_{18}x_{114} -x_{114} +x_{116} \\[3pt]
S_{101} &={}2 -4x_{1} +2x_{1}^{2} -2x_{1}x_{15} -2x_{1}x_{16} -2x_{1}x_{115} \\*
&\quad{}-2x_{1}x_{116} +2x_{15} +x_{15}^{2} +2x_{15}x_{115} +2x_{16} +x_{16}^{2} \\*
&\quad{}+2x_{16}x_{116} -x_{20}x_{117} +x_{103}x_{113} +2x_{115} +x_{115}^{2} +2x_{116} \\*
&\quad{}+x_{116}^{2} \\[3pt]
S_{102} &={}-x_{11}x_{20} -x_{12}x_{20} -x_{13}x_{20} -x_{14}x_{20} -x_{15}x_{20} -x_{16}x_{20} \\*
&\quad{}-x_{17}x_{20} -x_{18}x_{20} -x_{19}x_{20} -10x_{20} +x_{118} \\[3pt]
S_{103} &={}1 -x_{119} -x_{119}x_{166} -x_{119}x_{167} -x_{119}x_{217} +x_{172} \\*
&\quad{}+x_{173} +x_{179} +x_{184} +x_{189} +x_{220} \\[3pt]
S_{104} &={}-x_{20} -x_{20}x_{206} -x_{20}x_{207} -x_{20}x_{212} +x_{120} \\[3pt]
S_{105} &={}-80 -2x_{0} -2x_{0}x_{123} -x_{1} -x_{1}x_{123} -x_{2} \\*
&\quad{}-x_{2}x_{123} -2x_{3} -2x_{3}x_{123} -x_{4} -x_{4}x_{123} -x_{5} \\*
&\quad{}-x_{5}x_{123} -x_{6} -x_{6}x_{123} -x_{8} -x_{8}x_{123} -2x_{11} \\*
&\quad{}-2x_{11}x_{123} -14x_{12} -14x_{12}x_{123} -x_{13} -x_{13}x_{123} -x_{14} \\*
&\quad{}-x_{14}x_{123} -x_{15} -x_{15}x_{123} -x_{16} -x_{16}x_{123} -x_{17} \\*
&\quad{}-x_{17}x_{123} -x_{18} -x_{18}x_{123} -x_{19} -x_{19}x_{123} -x_{20} \\*
&\quad{}-x_{20}x_{123} -4x_{22} -4x_{22}x_{123} -2x_{25} -2x_{25}x_{123} -x_{29} \\*
&\quad{}-x_{29}x_{123} -x_{31} -x_{31}x_{123} -x_{33} -x_{33}x_{123} -x_{35} \\*
&\quad{}-x_{35}x_{123} -2x_{39} -2x_{39}x_{123} -4x_{41} -4x_{41}x_{123} -2x_{42} \\*
&\quad{}-2x_{42}x_{123} -x_{45} -x_{45}x_{123} -2x_{67} -2x_{67}x_{123} -x_{118} \\*
&\quad{}-x_{118}x_{123} -x_{119} -x_{119}x_{123} -80x_{123} -x_{123}x_{166} -x_{123}x_{167} \\*
&\quad{}-x_{123}x_{194} -x_{123}x_{195} -x_{123}x_{201} -x_{123}x_{206} -x_{123}x_{207} -x_{123}x_{212} \\*
&\quad{}-x_{123}x_{217} +x_{124} -x_{166} -x_{167} -x_{194} -x_{195} \\*
&\quad{}-x_{201} -x_{206} -x_{207} -x_{212} -x_{217} \\[3pt]
S_{106} &={}-2x_{0}x_{124} -x_{1}x_{124} -x_{2}x_{124} -2x_{3}x_{124} -x_{4}x_{124} -x_{5}x_{124} \\*
&\quad{}-x_{6}x_{124} -x_{8}x_{124} -2x_{11}x_{124} -14x_{12}x_{124} -x_{13}x_{124} -x_{14}x_{124} \\*
&\quad{}-x_{15}x_{124} -x_{16}x_{124} -x_{17}x_{124} -x_{18}x_{124} -x_{19}x_{124} -x_{20}x_{124} \\*
&\quad{}-4x_{22}x_{124} -2x_{25}x_{124} -x_{29}x_{124} -x_{31}x_{124} -x_{33}x_{124} -x_{35}x_{124} \\*
&\quad{}-2x_{39}x_{124} -4x_{41}x_{124} -2x_{42}x_{124} -x_{45}x_{124} -2x_{67}x_{124} -x_{118}x_{124} \\*
&\quad{}-x_{119}x_{124} -81x_{124} -x_{124}x_{166} -x_{124}x_{167} -x_{124}x_{194} -x_{124}x_{195} \\*
&\quad{}-x_{124}x_{201} -x_{124}x_{206} -x_{124}x_{207} -x_{124}x_{212} -x_{124}x_{217} +x_{125} \\[3pt]
S_{107} &={}-1 +x_{121}^{2} +x_{124}^{2} -x_{125}^{2} \\[3pt]
S_{108} &={}-2 -2x_{119} -x_{119}^{2} +2x_{119}x_{121} +2x_{119}x_{122} -x_{120}x_{126} \\*
&\quad{}+2x_{121} +2x_{122} \\[3pt]
S_{109} &={}-4x_{128} -4x_{128}x_{133} +x_{134} \\[3pt]
S_{110} &={}-x_{121}x_{128} -x_{122}x_{128} +x_{135} \\[3pt]
S_{111} &={}-1 +x_{127}^{2} +x_{128}^{2} -x_{135}^{2} \\[3pt]
S_{112} &={}-x_{121}x_{130} -x_{122}x_{130} +x_{136} \\[3pt]
S_{113} &={}-1 +x_{129}^{2} +x_{130}^{2} -x_{136}^{2} \\[3pt]
S_{114} &={}-x_{132} -x_{132}x_{134} +x_{137} \\[3pt]
S_{115} &={}-1 +x_{131}^{2} +x_{132}^{2} -x_{137}^{2} \\[3pt]
S_{116} &={}-x_{128}^{2} +x_{139} \\[3pt]
S_{117} &={}x_{130} -x_{138}x_{139} -x_{139} \\[3pt]
S_{118} &={}1 -x_{121} -x_{122} -x_{129}x_{140} +x_{134} \\[3pt]
S_{119} &={}-x_{127} -x_{129}x_{141} +x_{131} \\[3pt]
S_{120} &={}-25 -x_{11} -13x_{12} -4x_{22} -2x_{25} -x_{29} \\*
&\quad{}-x_{31} -x_{33} -x_{35} -2x_{39} -4x_{41} -2x_{42} \\*
&\quad{}-x_{45} -2x_{67} -4x_{128}x_{142} +x_{132} \\[3pt]
S_{121} &={}-x_{143}x_{144} +x_{145} \\[3pt]
S_{122} &={}-1 -2x_{121}x_{145} -2x_{122}x_{145} +x_{143}^{2} +x_{144}^{2} \\[3pt]
S_{123} &={}-2 -2x_{3} -x_{3}^{2} +2x_{3}x_{121} +2x_{3}x_{122} +2x_{121} \\*
&\quad{}+2x_{122} -x_{146} \\[3pt]
S_{124} &={}x_{3}x_{143} -x_{4} +x_{143} -x_{144} -x_{146}x_{147} \\[3pt]
S_{125} &={}-x_{148}x_{149} +x_{150} \\[3pt]
S_{126} &={}-1 -2x_{121}x_{150} -2x_{122}x_{150} +x_{148}^{2} +x_{149}^{2} \\[3pt]
S_{127} &={}-2 -2x_{5} -x_{5}^{2} -2x_{5}x_{6} +2x_{5}x_{121} +2x_{5}x_{122} \\*
&\quad{}-2x_{6} -x_{6}^{2} +2x_{6}x_{121} +2x_{6}x_{122} +2x_{121} +2x_{122} \\*
&\quad{}-x_{151} \\[3pt]
S_{128} &={}x_{5}x_{148} +x_{6}x_{148} -x_{20} +x_{148} -x_{149} -x_{151}x_{152} \\[3pt]
S_{129} &={}-1 -x_{118}^{2} +2x_{118}x_{121} +2x_{118}x_{122} -x_{153} \\[3pt]
S_{130} &={}x_{118}x_{148} -x_{119} -x_{149} -x_{153}x_{154} \\[3pt]
S_{131} &={}-1 -x_{119}^{2} +2x_{119}x_{121} +2x_{119}x_{122} -x_{155} \\[3pt]
S_{132} &={}-1 +x_{119}x_{143} -x_{144} -x_{155}x_{156} -x_{166} -x_{167} \\*
&\quad{}-x_{217} \\[3pt]
S_{133} &={}-x_{157}x_{158} +x_{159} \\[3pt]
S_{134} &={}-1 -2x_{121}x_{159} -2x_{122}x_{159} +x_{157}^{2} +x_{158}^{2} \\[3pt]
S_{135} &={}-1 +x_{119}x_{157} -x_{155}x_{160} -x_{158} -x_{194} -x_{195} \\*
&\quad{}-x_{201} \\[3pt]
S_{136} &={}-1 +x_{119}x_{148} -x_{149} -x_{155}x_{161} -x_{206} -x_{207} \\*
&\quad{}-x_{212} \\[3pt]
S_{137} &={}-x_{162}x_{163} +x_{164} \\[3pt]
S_{138} &={}-1 -2x_{121}x_{164} -2x_{122}x_{164} +x_{162}^{2} +x_{163}^{2} \\[3pt]
S_{139} &={}-1 -x_{5} -x_{5}x_{166} -x_{5}x_{167} -x_{5}x_{217} -x_{6} \\*
&\quad{}-x_{6}x_{166} -x_{6}x_{167} -x_{6}x_{217} +x_{119}x_{162} +x_{162} -x_{163} \\*
&\quad{}-x_{165} -x_{165}x_{172} -x_{165}x_{173} -x_{165}x_{179} -x_{165}x_{184} -x_{165}x_{189} \\*
&\quad{}-x_{165}x_{220} -2x_{166} -x_{167} -x_{217} \\[3pt]
S_{140} &={}-x_{168}x_{169} +x_{170} \\[3pt]
S_{141} &={}-1 -2x_{121}x_{170} -2x_{122}x_{170} +x_{168}^{2} +x_{169}^{2} \\[3pt]
S_{142} &={}-x_{20} -x_{20}x_{172} -x_{20}x_{173} -x_{20}x_{179} -x_{20}x_{184} -x_{20}x_{189} \\*
&\quad{}-x_{20}x_{220} +x_{174} \\[3pt]
S_{143} &={}x_{119}x_{168} +x_{168} -x_{169} -x_{171}x_{174} -x_{172} \\[3pt]
S_{144} &={}-x_{175}x_{176} +x_{177} \\[3pt]
S_{145} &={}-1 -2x_{121}x_{177} -2x_{122}x_{177} +x_{175}^{2} +x_{176}^{2} \\[3pt]
S_{146} &={}x_{119}x_{175} -x_{174}x_{178} +x_{175} -x_{176} -x_{179} \\[3pt]
S_{147} &={}-x_{180}x_{181} +x_{182} \\[3pt]
S_{148} &={}-1 -2x_{121}x_{182} -2x_{122}x_{182} +x_{180}^{2} +x_{181}^{2} \\[3pt]
S_{149} &={}x_{119}x_{180} -x_{174}x_{183} +x_{180} -x_{181} -x_{184} \\[3pt]
S_{150} &={}-x_{185}x_{186} +x_{187} \\[3pt]
S_{151} &={}-1 -2x_{121}x_{187} -2x_{122}x_{187} +x_{185}^{2} +x_{186}^{2} \\[3pt]
S_{152} &={}x_{119}x_{185} -x_{174}x_{188} +x_{185} -x_{186} -x_{189} \\[3pt]
S_{153} &={}-x_{190}x_{191} +x_{192} \\[3pt]
S_{154} &={}-1 -2x_{121}x_{192} -2x_{122}x_{192} +x_{190}^{2} +x_{191}^{2} \\[3pt]
S_{155} &={}-x_{20} -x_{20}x_{194} -x_{20}x_{195} -x_{20}x_{201} +x_{196} \\[3pt]
S_{156} &={}x_{119}x_{190} +x_{190} -x_{191} -x_{193}x_{196} -x_{194} \\[3pt]
S_{157} &={}-x_{197}x_{198} +x_{199} \\[3pt]
S_{158} &={}-1 -2x_{121}x_{199} -2x_{122}x_{199} +x_{197}^{2} +x_{198}^{2} \\[3pt]
S_{159} &={}x_{119}x_{197} -x_{196}x_{200} +x_{197} -x_{198} -x_{201} \\[3pt]
S_{160} &={}-x_{202}x_{203} +x_{204} \\[3pt]
S_{161} &={}-1 -2x_{121}x_{204} -2x_{122}x_{204} +x_{202}^{2} +x_{203}^{2} \\[3pt]
S_{162} &={}x_{119}x_{202} -x_{120}x_{205} +x_{202} -x_{203} -x_{206} \\[3pt]
S_{163} &={}-x_{208}x_{209} +x_{210} \\[3pt]
S_{164} &={}-1 -2x_{121}x_{210} -2x_{122}x_{210} +x_{208}^{2} +x_{209}^{2} \\[3pt]
S_{165} &={}x_{119}x_{208} -x_{120}x_{211} +x_{208} -x_{209} -x_{212} \\[3pt]
S_{166} &={}-x_{213}x_{214} +x_{215} \\[3pt]
S_{167} &={}-1 -2x_{121}x_{215} -2x_{122}x_{215} +x_{213}^{2} +x_{214}^{2} \\[3pt]
S_{168} &={}-x_{20} -x_{20}x_{166} -x_{20}x_{167} -x_{20}x_{217} +x_{218} \\[3pt]
S_{169} &={}x_{119}x_{213} +x_{213} -x_{214} -x_{216}x_{218} -x_{217} \\[3pt]
S_{170} &={}x_{119}x_{128} +x_{127} +x_{128} -x_{135} -x_{174}x_{219} -x_{220} \\[3pt]
S_{171} &={}-326 -48x_{0} -2x_{0}^{2} -2x_{0}x_{3} +2x_{0}x_{143} +2x_{0}x_{162} \\*
&\quad{}-x_{1}^{2} +2x_{1}x_{148} -4x_{2} -x_{2}^{2} -2x_{2}x_{8} +2x_{2}x_{157} \\*
&\quad{}-24x_{3} -x_{3}^{2} +2x_{3}x_{162} -4x_{8} -x_{8}^{2} +2x_{8}x_{157} \\*
&\quad{}-3x_{11} -x_{11}^{2} +2x_{11}x_{168} -15x_{12} -x_{12}^{2} +2x_{12}x_{175} \\*
&\quad{}-2x_{13} -x_{13}^{2} +2x_{13}x_{180} -2x_{14} -x_{14}^{2} +2x_{14}x_{185} \\*
&\quad{}-2x_{15} -x_{15}^{2} +2x_{15}x_{190} -2x_{16} -x_{16}^{2} +2x_{16}x_{197} \\*
&\quad{}-2x_{17} -x_{17}^{2} +2x_{17}x_{202} -2x_{18} -x_{18}^{2} +2x_{18}x_{208} \\*
&\quad{}-2x_{19} -x_{19}^{2} +2x_{19}x_{213} -4x_{22} -2x_{25} -x_{29} \\*
&\quad{}-x_{31} -x_{33} -x_{35} -2x_{39} -4x_{41} -2x_{42} \\*
&\quad{}-x_{45} -2x_{67} -x_{121}x_{221} -x_{122}x_{221} +x_{128} +24x_{143} \\*
&\quad{}-x_{143}^{2} -x_{148}^{2} +4x_{157} -x_{157}^{2} +24x_{162} -x_{162}^{2} \\*
&\quad{}+2x_{168} -x_{168}^{2} +2x_{175} -x_{175}^{2} +2x_{180} -x_{180}^{2} \\*
&\quad{}+2x_{185} -x_{185}^{2} +2x_{190} -x_{190}^{2} +2x_{197} -x_{197}^{2} \\*
&\quad{}+2x_{202} -x_{202}^{2} +2x_{208} -x_{208}^{2} +2x_{213} -x_{213}^{2} \\*
&\quad{}+x_{221}
\end{align*}
\endgroup


\section{The four fixed integers in the compact equations}
\label{sec:compact-constants}

This appendix specifies the complete coefficient tuple used in
Section~\ref{sec:compact}.  It is the indexed BFS943 \emph{radix-67}
tuple of the retained source, including its corrected initializer and
the single source shift in \(M_F\).\cite{indexed,certificate}
No radix-55 coefficient is substituted.  All indices, finite sets,
digits, sums, floors, and decoding operations in this appendix define
constants; they introduce no quantified unknowns into the equation.

\subsection{Finite formulas for the coefficient integers}

Put
\[
 \rho=2^{67},\qquad
 \mathcal B=\rho^{8807},\qquad
 \beta=\mathcal B-1,\qquad
 G_m(x)=\sum_{j=0}^{m-1}x^j.
\]
The symbol \(\mathcal B\) here is a fixed radix power, distinct from the
polynomial abbreviation \(B=(K_0+w)C-F\) in the compact equation.
For integers \(a\le b\), write
\([a,b]_{\mathbb Z}=\{a,a+1,\ldots,b\}\), and let an interval with
upper endpoint smaller than its lower endpoint be empty.  Define
\begin{align*}
 \mathcal E={}&[0,942]_{\mathbb Z}
   \cup[983,992]_{\mathbb Z}\cup[1003,1011]_{\mathbb Z}\\
 &{}\cup[1013,1031]_{\mathbb Z}
   \cup\{1063,1064,1095,2195,3296\}.
\end{align*}
Then
\begin{align}
 M_C={}&\beta-\sum_{e\in\mathcal E\setminus\{1\}}\rho^e
                    -2\rho^{1063},\nonumber\\
 M_F={}&\beta+4+(2^{51}-2)(\rho^{4393}+\rho^{5495})\nonumber\\
 &+\rho^{1013}G_{10}(\rho)+\rho^{1065}G_9(\rho)
          +\rho^{2195}+\rho^{3296}\nonumber\\
 &+(2^{65}-1)\rho^{8792}G_{13}(\rho)
          +(2^{19}-1)\rho^{8805},                 \label{eq:compact-fixedmasks}\\
 K_0={}&\rho^{30}+\rho^{62}+\rho^{1100}
       +\rho^{2201}+\rho^{8849}
       +(\rho^{4393}+\rho^{5495})\Phi(1/\rho).
                                                     \label{eq:compact-fixedK}
\end{align}
The polynomial \(\Phi\) is the following explicitly specified finite
polynomial:
\begin{align}
 \Phi(x)={}&(1+2^{51})G_{943}(x)
       -2^{10}\sum_{i=0}^{942}v(t_i)x^i\nonumber\\
 &+\sum_{(b,e,m)\in\mathcal C}2^b x^eG_m(2x),
                                                     \label{eq:compact-fixedPhi}\\
 \mathcal C={}&\{
   (10,1023,8),(18,1013,9),(27,1003,8),(35,983,9),\nonumber\\
 &\hspace{8mm}(44,1031,1),(45,1022,1),(46,1011,1),(47,992,1),\nonumber\\
 &\hspace{8mm}(48,1095,1),(49,2195,1),(50,3296,1)\}.
                                                     \label{eq:compact-fixedcontrols}
\end{align}
The tiles \(t_i\) and their integer encoding \(v\) are given below.
Every exponent occurring in \(\Phi\) is at most \(3296<4393\);
therefore the reciprocal substitution in
\eqref{eq:compact-fixedK} cancels exactly, leaving an integer.
Equations~\eqref{eq:compact-fixedmasks}--\eqref{eq:compact-fixedcontrols}
are fixed finite recipes, not arithmetic involving unknown exponents.

\subsection{The 943 fixed indexed tiles}

First decode the 430 triples \((u_j,v_j,d_j)\) by the next subsection.
For a state index \(r\) and bit \(b\in\{0,1\}\), set
\[
 \mathfrak h(r,b)=2(r+1)+b.
\]
Let
\[
 \mathcal R=\{u_j:d_j=1,\ 0\le j<430\},\qquad
 \mathcal L=\{u_j:d_j=0,\ 0\le j<430\},
\]
and, for each state index \(r\), define
\[
 \mathsf r_r=6+\bigl|\mathcal R\cap[0,r-1]_{\mathbb Z}\bigr|,
 \qquad
 \mathsf l_r=205+\bigl|\mathcal L\cap[0,r-1]_{\mathbb Z}\bigr|.
\]
The notation \(\mathsf r_r\) and \(\mathsf l_r\) denotes two fixed
integer-valued tables.  The following table specifies every tile
\(t_i=(w,n,e,s)\), \(0\le i\le942\), in its exact indexed order.

\begingroup
\small
\setlength{\tabcolsep}{5pt}
\begin{longtable}{@{}ll@{}}
\toprule
Tile index \(i\) & Tile \(t_i=(w,n,e,s)\)\\
\midrule
\endfirsthead
\toprule
Tile index \(i\) & Tile \(t_i=(w,n,e,s)\)\\
\midrule
\endhead
\bottomrule
\endfoot
\(0\) & \((2,435,3,1)\)\\
\(1\) & \((5,0,0,0)\)\\
\(j+2,\quad j=2r+b<430\)
 & \(\bigl((1-d_j)\mathsf l_{u_j},\mathfrak h(r,b),
                 d_j\mathsf r_{u_j},v_j\bigr)\)\\
\(432+b,\quad b=0,1\)
 & \(\bigl(0,\mathfrak h(215,b),0,\mathfrak h(215,b)\bigr)\)\\
\(434+b,\quad b=0,1\) & \((0,b,0,b)\)\\
\(424+2\mathsf r_r+b,\quad r\in\mathcal R,\ b=0,1\)
 & \(\bigl(\mathsf r_r,b,0,\mathfrak h(r,b)\bigr)\)\\
\(424+2\mathsf l_r+b,\quad r\in\mathcal L,\ b=0,1\)
 & \(\bigl(0,b,\mathsf l_r,\mathfrak h(r,b)\bigr)\)\\
\(924+j,\quad 0\le j<4\)
 & \(\bigl(0,j+430\lfloor j/2\rfloor,0,438\bigr)\)\\
\(928\) & \((2,434,2,0)\)\\
\(929\) & \((3,436,4,4)\)\\
\(930\) & \((4,436,4,0)\)\\
\(497+t,\quad t=434,435,436\) & \((1,438,1,t)\)\\
\(\substack{934+3a+b-6\lfloor a/4\rfloor\\
             a\in\{0,1,4\},\ b=0,1,2}\)
 & \((a,437,b,437)\)\\
\end{longtable}
\endgroup

For a fixed tile \(t=(w,n,e,s)\), the encoding used in
\eqref{eq:compact-fixedPhi} is
\[
 v(t)=w+2^8n+2^{17}e+2^{25}s+127\cdot2^{34}.
\]
These tile coordinates are bound variables in a constant definition;
they are unrelated to the independent positive unknowns that happen to
use some of the same letters.

\subsection{A closed decoder for all 430 actions}

For nonnegative integers \(A,k\) and an integer base \(b\ge2\), define
\[
 D_b(A,k)=\left\lfloor\frac{A}{b^k}\right\rfloor\bmod b,
 \qquad
 a\bmod b=a-b\left\lfloor\frac ab\right\rfloor.
\]
Let \(\mathcal A\) be the fixed integer specified in the next subsection.
Set
\begin{align*}
 N&=\mathcal A\bmod2^{636},\\
 T_0&=\left\lfloor\mathcal A/2^{636}\right\rfloor\bmod2^{1675},\\
 E_0&=\left\lfloor\mathcal A/2^{2311}\right\rfloor.
\end{align*}
For \(0\le r<215\), \(b\in\{0,1\}\), and \(j=2r+b\), put
\begin{align}
 \nu_j&=D_2(N,j),&
 \sigma_j&=\sum_{k=0}^{j-1}\nu_k,\nonumber\\
 m_r&=D_2(N,430+r),&
 \tau_r&=\sum_{k=0}^{r-1}m_k,\nonumber\\
 F_r&=14(1-m_r)+m_rD_{16}(E_0,\tau_r),\nonumber\\
 u_j&=\nu_j(1+\sigma_j)
          +(1-\nu_j)D_{216}(T_0,j-\sigma_j),\nonumber\\
 v_j&=D_2(F_r,2b),&
 d_j&=D_2(F_r,2b+1).                              \label{eq:compact-actiondecoder}
\end{align}
Empty sums are zero.  In particular,
\((v_0,d_0,u_0)=(1,1,1)\).  The constructor uses this actual first
target, including the source's state renaming.

\subsection{The complete 436-character seed}

The alphabet in increasing digit order \(0,1,\ldots,61\) is
\begin{center}
\texttt{0123456789ABCDEFGHIJKLMNOPQRSTUVWXYZabcdefghijklmnopqrstuvwxyz}.
\end{center}
Concatenate the following seven full lines and final two-character line
with no spaces or newlines.  If the resulting 436 digits have values
\(e_0,\ldots,e_{435}\), define
\[
 \mathcal A=\sum_{j=0}^{435}e_j62^{435-j}.
\]
\begin{lstlisting}[language={},basicstyle=\ttfamily\footnotesize,
  columns=fixed,keepspaces=true,breaklines=false]
BxEnhvy6NipzkWmROo0GsA3uK3twWLxscLq20UTumy5Cdu461LwD6g7RDNzF8Y
iYDFTpB8ehAypd7d6OsOL4APamnJeNXQypUezCmW9n4oXP6XfgmiWuAWZ8ZJ7U
OP7gu4kVdlGZqZHULQMar1dXNeJW4i7oNVI8usp8BlovjjZbwPCovZI328fRUq
4q92R81s5GjWQaQfGhkq9UYMK3rRMaBmsOIBkJqc2cj2VRIZ6mjzkBIOkW9SjZ
FzeA6J5u4j3sBL3SdgwXzM9mHvQivu8EJ1TBTpRhHDqZIbmYgwQ8VqTAFotl6F
We1YC554GKrVHsA8oUGKSA6F6k5KREBSD0Tgd7bFALY84LLbpw574BQyvdoHrl
0Oe35uwDf15DpT6VJWZBjSbHPpPAINYTte0eUxnTFTKBlGj8Am6sb1gLldh8WZ
8P
\end{lstlisting}

\subsection{Literal bytes and reconstruction check}

The four resulting integers have the following exact bit lengths:
\begin{center}
\begin{tabular}{@{}lr@{}}
\toprule
Fixed integer & Bit length\\
\midrule
\(\beta\) & 590\,069\\
\(M_C\) & 590\,069\\
\(M_F\) & 590\,070\\
\(K_0\) & 592\,884\\
\bottomrule
\end{tabular}
\end{center}
The artifact contains their canonical unsigned little-endian bytes,
compressed separately, with both compressed and raw SHA-256 identities
in each graph's coefficient binding.  The included
\texttt{rebuild\_coefficients.py} independently reconstructs the decoded
actions, all 943 indexed tiles, and all four integers from this finite
recipe.  It verifies every raw byte; \texttt{verify.py} also verifies
the inequalities \eqref{eq:compact-portbox}.  The table and seed therefore give a complete
finite specification; the coefficient files provide a convenient second
representation of the same four constants, with no remaining compiler
step needed to define the polynomial.


\section{File map and reproduction}
\begin{description}[style=nextline]
\item[\code{report.tex}, \code{report_inputs/}, \code{report.pdf}]
This report, its editable sections, and the compiled PDF.
\item[\code{arithmetic/}]
The selected unrestricted-integer and natural-domain quartics, all collected coefficients, source mutation certificates, restoration maps, exact checkers, and retained ordinary proof dependencies.
\item[\code{compact/}]
Four complete ten-positive circuits, four fixed coefficient payloads, deterministic reconstruction, full new proof, inherited proof sources, and separate independent checker.
\item[\code{reviews/}, \code{verification/}]
Mathematical review notes and actual verification receipts, with their scopes explicitly stated.
\item[\code{PROVENANCE.json}, \code{MANIFEST.sha256}]
Immutable comparison and source identities, plus the exact release-file checksum list.
\end{description}
From the extracted package, run
\begin{lstlisting}
python3 verify_all.py
pdflatex -interaction=nonstopmode -halt-on-error report.tex
pdflatex -interaction=nonstopmode -halt-on-error report.tex
\end{lstlisting}
Verification requires Python's standard library and performs no network access. Rebuilding the PDF requires a standard LaTeX installation. PDF metadata can differ between TeX distributions without changing any polynomial.

\begin{thebibliography}{20}
\bibitem{repository} Vladimir Reshetnikov and contributors. \emph{ProveIt, DiophantineConsistency}. Comparison commit \code{e0c98c629bf10d30b695884f6b1d2ab361bdf513}, 25 October 2026. \url{https://github.com/VladimirReshetnikov/ProveIt/tree/e0c98c629bf10d30b695884f6b1d2ab361bdf513/Computability/DiophantineConsistency}.
\bibitem{cofinal} \emph{Cofinal anchors and height one equations for ZFC consistency}, Round 17, 25 October 2026. Incoming archive \code{ConZFC_Cofinal_Anchors_Height_One_Round17_2026-10-25.zip}, in the pinned repository. In particular \code{ORIENTATION_FIRST_PROOF.md}, \code{STRICT_BUDGET_PROOF.md}, and \code{ONE_ROW_SEMANTIC_ANCHOR_PROOF.md}.
\bibitem{graphdigits} \emph{Smaller explicit Diophantine equations for Con(ZFC)}, 24 October 2026. Incoming archive \code{conzfc_graph_digits_2026-10-24.zip}. The four-group natural low-digit pooling construction.
\bibitem{certificate} \emph{ConZFC Certificate Reductions}, 25 October 2026. Incoming archive \code{ConZFC_Certificate_Reductions_2026-10-25.zip}. The $78/630$ and $79/468$ scale-free compact parents and their full coefficient recipe.
\bibitem{unshift} ProveIt. \emph{Ten-positive unshifted transport}. \code{research/ten-positive-unshift-t-square}, \code{research/ten-positive-unshift-t-bv}, and \code{research/ten-positive-unshift-t-finite}, at the comparison commit. The complete ordinary proofs are copied in \code{compact/proof_sources/}.
\bibitem{indexed} ProveIt. \emph{Indexed tile enumeration}. The full coherent radix-$2^{67}$ BFS943 coefficient recipe and source calibration, retained in \code{compact/proof_sources/indexed_native.md} and \code{indexed_constants.tex}.
\bibitem{native61} ProveIt contribution 61. \emph{Congruence carriers}, with later native, saturation, ratio, and product-square proofs. The relevant ordinary Pell and reconstruction proofs are retained in \code{compact/proof_sources/}, with immutable identities in \code{compact/SOURCE_LEDGER.json}.
\bibitem{endpoint} ProveIt. \code{SetTheory.BusyBeaver.ZFC.KEndpoint.c222_iff_zf_refutation}. The machine-independent numerical-calculus calibration. Source statement and proof are retained with the direct source dependencies.
\bibitem{choice} ProveIt. \code{SetTheory.BusyBeaver.ZFC.ConstructibleChoice.zf_refutation_iff_zfc_refutation}. Constructible-model removal of Choice. Source statement and proof are retained with the direct source dependencies.
\bibitem{sun} Zhi-Wei Sun. \emph{Further results on Hilbert's Tenth Problem}. Science China Mathematics 64 (2021), 281--306. The nonnegative-integer representation $a^2+b^2+c(c+1)$ is also used throughout the inherited direct arithmetic. \url{https://arxiv.org/abs/1704.03504}.
\end{thebibliography}
\end{document}
